\documentclass[11pt]{article}

\usepackage[a4paper,margin=1in]{geometry}
\usepackage{amsmath,amssymb,amsthm,mathtools}
\usepackage{microtype}
\usepackage[hidelinks]{hyperref}

\newtheorem{theorem}{Theorem}[section]
\newtheorem{lemma}[theorem]{Lemma}
\newtheorem{proposition}[theorem]{Proposition}
\newtheorem{corollary}[theorem]{Corollary}
\theoremstyle{definition}
\newtheorem{definition}[theorem]{Definition}
\theoremstyle{remark}
\newtheorem{remark}[theorem]{Remark}

\newcommand{\cA}{\mathcal A}
\newcommand{\cB}{\mathcal B}
\newcommand{\cC}{\mathcal C}
\newcommand{\cY}{\mathcal Y}
\newcommand{\R}{\mathbb R}
\newcommand{\Z}{\mathbb Z}
\newcommand{\F}{\mathbb F}
\newcommand{\1}{\mathbf 1}
\newcommand{\Cov}{\operatorname{Cov}}
\newcommand{\Var}{\operatorname{Var}}
\newcommand{\rank}{\operatorname{rank}}
\newcommand{\rect}{\operatorname{rect}}
\newcommand{\dist}{\operatorname{dist}}
\newcommand{\Dcc}{\operatorname{D}}

\title{Logarithmic Communication for Low-Degree Slices and Product Domains\\
with Spectral--Lattice Stability}
\author{Anonymous draft}
\date{August 26, 2026}

\begin{document}
\maketitle

\begin{abstract}
We study the adjacent two-value cross-intersection problem that occurs in an
equivalent formulation of the Log-Rank Conjecture.  We prove a
spectral--lattice rectangle theorem for an arbitrary sparse right family: a
covariance gap away from a real defect space controls the Hamming distance of
every opposite row from a Boolean template, while low covariance on an
integer generating set yields a large joint atom by a discrete entropy
argument.  The resulting rectangle bound depends on an exact Boolean tube
radius and on the entropy of integer generator statistics, but has no
exponential dependence on the dimension of the defect space or on its number
of Boolean points.

We also prove an atomic product transfer theorem.  If every Boolean
degree-one function on each factor is an atom or its complement, and atom
families have polynomial growth in centered rank with exponent $\alpha$, then
every row-degree-one Boolean matrix on the Cartesian product has deterministic
communication complexity at most $\alpha\log_2(\rank M+1)+O(1)$.  Conversely,
a degree-one selector partition embeds arbitrary Boolean matrices, giving an
abstract tractable/universal dichotomy.

In a different direction, let a Boolean matrix have the complete Johnson
slice $\binom{[n]}k$ as its column domain, and suppose every row has slice
degree at most $s$.  Combining slice query-complexity bounds with rank pooling
and a new symmetric-tree canonicalization, we prove
\[
 \Dcc(M)=O\bigl(s^8\log(\rank M+1)+s^{16}\bigr)
\]
whenever $\min(k,n-k)\ge144s^4$.  Thus the
Log-Rank Conjecture holds throughout this class even when
$s=\operatorname{polylog}(\rank M)$.  The same bound holds on the full cube
without a layer-depth hypothesis, with a sharper dependence on $s$; in the
deeper regime $\min(k,n-k)\ge C^s$, we also obtain the sharper bound
$O(s^3(\log(\rank M+1)+s))$.  At the opposite boundary $k\le s$, every
Boolean function on the layer already has degree at most $s$, so arbitrary
Boolean matrices occur.  Qualitatively, the deep-layer hypothesis can be
replaced by the sharp junta threshold $\min(k,n-k)\ge2s$: for every fixed
$s$, we obtain $\Dcc(M)=O_s(\log(\rank M+1))$.
More generally, if every row on a symmetric or complete-product coordinate
domain has bounded-answer query depth $q$, then
\[
 \Dcc(M)\le q\bigl(\lceil\log_2(\rank M+1)\rceil+
 (q+1)\lceil\log_2Q\rceil\bigr),
\]
where $Q$ bounds the number of answers to a query.
The same bounded-answer transfer proves
$O_Q(s^8\log(\rank M+1)+s^{16})$ communication for degree-$s$ rows on
complete product alphabets of size at most $Q$, and on $Q$-color multislices
whose minimum color multiplicity is at least $216s^4$.
An unisolvent-restriction theorem extends all these protocols to sparse
column subsets.  In particular, a uniform random
$O(\binom ns\log\binom ns)$-column subfamily of a Johnson layer is, with high
probability, simultaneously sufficient for every degree-$s$ Boolean row
family.  Projection-DPP volume sampling removes the logarithmic factor and
selects exactly the information-theoretic minimum $\binom ns$ columns with
probability one; for fixed $s$ in a middle layer this is exponentially sparse.
The same minimal projection-DPP construction applies to bounded-alphabet
product degree spaces, whose dimension is $O_{Q,s}(n^s)$ against an ambient
domain of size at least $2^n$.
At degree one, twice-degree unisolvence removes any ambient-extension
assumption: every fixed degree-two-unisolvent Johnson column family
simultaneously satisfies the adjacent-two-value theorem for every row family.
Such a family exists with exactly $\binom n2$ columns, for example in the
explicit radius-two cross of Theorem~\ref{thm:explicit-cross}; projection-DPP
sampling gives an invariant alternative.  It satisfies
\[
 \Dcc(M)\le\lceil\log_2(\rank M+1)\rceil+1,
 \qquad \rect(M)\ge1/[4(\rank M+1)].
\]
On every fixed balanced range $\beta n\le k\le(1-\beta)n$, a
quadratic-scale hypergeometric two-point escape theorem reduces the universal
column count further to $O_\beta(n)$, with the same bounds.  A strengthened construction needs no
protected rank anchor: the same random column family is simultaneously a
two-value detector and a robust degree-one interpolation set.  It retains the
theorem after adversarial deletion of a $\beta^2/10$-fraction of \emph{all}
its columns.
The underlying hypergeometric escape constant is exactly $2/5$ after
quadratic normalization, with unique equality at the central $J(6,3)$
three-marked distribution.
There is also a deterministic spectral version: any $k$-uniform column
family whose empirical incidence covariance has eigenvalue greater than
$n/[8(n-2)]$ on $\mathbf1^\perp$ satisfies the same logarithmic protocol.
Balanced $2$-designs meet this condition in a broad parameter range; the
Sylvester--Hadamard construction gives an explicit family of exactly $n$
columns for every $n=2^t-1$, $t\ge3$.

As a concrete application, partition $[d]$ into arbitrary blocks and let the
right family choose a fixed number $k_i$ from each block, where
$2\le k_i\le n_i-2$.  This deterministic family can have exponentially small
relative density in the global layer and is generally not a complete
affine-code section.  Nevertheless every adjacent two-value opposite family
has the exact row classification: every nonconstant row is a coordinate
function or its complement.  This gives a monochromatic rectangle of relative
density at least $1/[4(d+1)]$ and the rank-sensitive communication bound
$\Dcc(M)\le\lceil\log_2(\rank M+1)\rceil+4$, which is optimal up to an
additive constant.
We also prove a robust transfer from a nearby rational subspace, with explicit
projector, Boolean-margin, and integer-height losses.  For layers at the
boundary $k_i=1$ or $n_i-1$, arbitrary Boolean matrices can
be embedded, so the interior hypothesis marks a sharp barrier.  More general
blockwise Johnson layers admit a bound governed by an explicit weighted
Boolean tube radius.  We also prove a logarithmic rank-sensitive protocol for
products of interior multislices over a bounded alphabet, with the same sharp
count-one boundary.  These are restricted-column results and do not solve the
general Log-Rank Conjecture.
\end{abstract}

\section{Introduction}

The Log-Rank Conjecture asks whether the deterministic communication
complexity of a Boolean matrix is polynomially bounded by the logarithm of
its real rank.  Equivalently, it predicts a quasipolynomial relation between
real rank and the partition number of a Boolean matrix.  The strongest known
general bound gives a monochromatic rectangle of relative density
$2^{-O(\sqrt r)}$ in every rank-$r$ Boolean matrix
\cite{SudakovTomon,HambardzumyanLovettShirley}.

Hambardzumyan, Lovett, and Shirley recently showed that Log-Rank is
equivalent to a relative rectangle conjecture for two set families whose
cross-intersections take two adjacent values
\cite{HambardzumyanLovettShirley}.  In that formulation one has
families $\cA,\cB\subseteq 2^{[d]}$ and an integer $a$ such that
\begin{equation}\label{eq:two-level}
 |A\cap B|\in\{a,a+1\}
 \qquad(A\in\cA,\ B\in\cB).
\end{equation}
The Boolean matrix
\begin{equation}\label{eq:matrix}
 M_{A,B}=\1_{\{|A\cap B|=a+1\}}
\end{equation}
has real rank at most $d+1$.  The conjecture asks for a monochromatic
rectangle retaining a $2^{-\operatorname{polylog}d}$ fraction of all
cross-pairs.

The problem remains difficult precisely when both families are sparse.
Classical and modern anti-concentration theorems control dense cube families
or independent coordinates \cite{Erdos,RaoYehudayoff,TaoVu,NguyenVu}; the
binary-scalar-product theorems of Kupavskii--Weltge and
Kupavskii--Tsarev control absolute products of full-span configurations
\cite{KupavskiiWeltge,KupavskiiTsarev}.  Neither type of conclusion is a
relative rectangle theorem for an arbitrarily sparse prescribed column
family.

Filmus--Ihringer classify Boolean degree-one functions on an individual
interior Johnson scheme as constants, dictators, and their complements
\cite{FilmusIhringer}.  Our local interval argument recovers this set-case
fact.  The contributions here are the product assembly, the cross-intersection
interpretation, the rank-sensitive protocol and sharp boundary universality,
and the spectral--lattice extension.

Our first contribution is a general sufficient condition formulated through
the covariance of the right family and the arithmetic geometry of a
low-mode subspace.  It is useful even when that subspace has dimension
linear in $d$.  Our second contribution is a particularly clean natural
application.  A third contribution is a general product transfer principle
that converts factorwise degree-one classifications and atom-rank growth into
Log-Rank protocols.  A fourth is a rank-pooling principle for low-degree
Boolean rows: a row-space basis compresses the common coordinate universe,
after which low-depth query trees become rank-sensitive communication
protocols.  A fifth is an anchor-free sparse transfer: one linear-size sample
simultaneously detects every forbidden two-value row and remains degree-one
unisolvent after adversarial column deletion.  A sixth replaces randomness
entirely by a deterministic empirical covariance gap and includes explicit
linear-size Hadamard-design columns.

\begin{theorem}[Interior product-Johnson theorem; informal]\label{thm:intro-product}
Partition $[d]$ into blocks of sizes $n_i$, and let $\cY$ consist of all sets
containing exactly $k_i$ points from block $i$, where
$2\le k_i\le n_i-2$.  If an arbitrary family $\cA$ satisfies
\eqref{eq:two-level} with $\cY$, then the
matrix \eqref{eq:matrix} satisfies
\[
 \rank_{\R}M\le d+1,
 \qquad
 \rect(M)\ge\frac1{4(d+1)},
 \qquad
 \Dcc(M)\le\lceil\log_2(\rank M+1)\rceil+4.
\]
\end{theorem}

\begin{theorem}[Anchor-free sparse HLS; informal]
\label{thm:intro-anchor-free}
Fix $0<\beta\le1/2$.  On every layer
$\beta n\le k\le(1-\beta)n$, a uniform random family of
$O_\beta(n+\log(1/\delta))$ columns has, with probability at least
$1-\delta$, the following simultaneous property.  After arbitrary deletion
of at most a $\beta^2/10$-fraction of all sampled columns, every
adjacent-two-value set-row matrix on the retained family satisfies
\[
 \Dcc(M)\le\lceil\log_2(\rank_{\R}M+1)\rceil+1,
 \qquad
 \rect(M)\ge\frac1{4(\rank_{\R}M+1)}.
\]
No subset of columns is designated as a protected rank anchor.
\end{theorem}

\begin{theorem}[Deterministic spectral sparse HLS; informal]
\label{thm:intro-spectral-sparse}
Let $Y\subseteq\binom{[n]}k$ be nonempty and let $\Sigma_Y$ be the empirical
covariance matrix of its incidence vectors.  If, for every
$v\perp\mathbf1$,
\[
 v^\top\Sigma_Yv\ge\lambda\|v\|_2^2,
 \qquad
 \lambda>\frac{n}{8(n-2)},
\]
then every adjacent-two-value set-row matrix on $Y$ satisfies
\[
 \Dcc(M)\le\lceil\log_2(\rank_{\R}M+1)\rceil+1,
 \qquad
 \rect(M)\ge\frac1{4(\rank_{\R}M+1)}.
\]
This is a deterministic condition on the prescribed column family.  In
particular, explicit Sylvester--Hadamard designs give $|Y|=n$ on the infinite
sequence $n=2^t-1$.
\end{theorem}

\begin{theorem}[Low-degree slice theorem; informal]
\label{thm:intro-low-degree}
Let $M$ be a Boolean matrix whose columns are indexed by $\binom{[n]}k$, and
suppose every row has slice degree at most $s\ge1$.  If
$\min(k,n-k)\ge144s^4$ and $r=\rank_{\R}M$, then, with $q=5184s^8$,
\[
 \Dcc(M)\le
 q\bigl(\lceil\log_2(r+1)\rceil+q+1\bigr).
\]
In particular, the Log-Rank Conjecture holds for this class whenever
$s=\operatorname{polylog}(r+1)$.
For each fixed $s$, the qualitative conclusion
$\Dcc(M)=O_s(\log(r+1))$ already holds under the sharp junta condition
$\min(k,n-k)\ge2s$.
\end{theorem}

The interior condition is sharp in a complexity-theoretic sense.  If a block
with $k_i=1$ is allowed, its columns are singletons and arbitrary Boolean
matrices can be encoded by taking each row to be its support.  Thus extending
the logarithmic theorem to the boundary would include the general Log-Rank
Conjecture.

Using the Filmus--Ihringer classification on multislices, the same product
and rank method gives a logarithmic protocol for products of multislices in
which every color has multiplicity at least two.  A color of multiplicity one
again embeds arbitrary Boolean matrices.

For fixed block parameters and a growing number of blocks, $\cY$ has exponentially
small relative density in the global middle layer.  It is deterministic and
is generally not the intersection of that layer with a binary affine coset.
Thus Theorem~\ref{thm:intro-product} is not a dense-subset, random-thinning,
or complete-affine-section corollary.

The proof has two complementary descriptions.  An elementary blockwise
hypergeometric variance computation shows that every row is within one
coordinate of a union of whole blocks.  More conceptually, the covariance
vanishes on the span of block indicators and has a constant gap on its
orthogonal complement; the Boolean tube at the two-point variance threshold
has Hamming radius one.

Our general theorem replaces the block span by an arbitrary integer-generated
real subspace.  The main parameters are an exact Boolean tube radius and the
entropy of the integer generator statistics.  A perturbative version permits
the true low-mode space to be merely close to the rational one.

\paragraph{Claim boundary.}
All results in this paper are restricted theorems.  They do not improve the
best known bound for every rank-$r$ Boolean matrix.  The product-Johnson class
does admit a full logarithmic protocol; the general spectral--lattice theorem
is only a one-rectangle result.  The publication-level novelty comparison is
still subject to independent expert verification.

\section{Two-level matrices and rectangle density}

For nonempty finite sets $X,Y$ and $M\in\{0,1\}^{X\times Y}$, define
\[
 \rect(M)=\max_{X_0\times Y_0\text{ monochromatic}}
 \frac{|X_0|}{|X|}\frac{|Y_0|}{|Y|}.
\]
Write $\Dcc(M)$ for deterministic communication complexity.  We use the
one-party-output convention; requiring both parties to output changes all
bounds below by at most one bit.

Under \eqref{eq:two-level},
\begin{equation}\label{eq:rank-factorization}
 M_{A,B}=|A\cap B|-a
 =-a+\sum_{j=1}^d\1_{\{j\in A\}}\1_{\{j\in B\}}.
\end{equation}
Hence
\begin{equation}\label{eq:rank-bound}
 \rank_{\R}M\le d+1.
\end{equation}

We repeatedly use the following elementary fact.

\begin{lemma}[Two-point variance]\label{lem:two-point-var}
If a real random variable is supported on two consecutive integers, then its
variance is at most $1/4$.
\end{lemma}

\begin{proof}
After translation it is Bernoulli with parameter $p$, and its variance is
$p(1-p)\le1/4$.
\end{proof}

\section{An atomic product transfer theorem}

We isolate the communication argument from the particular combinatorial
domains used later.  Let $\Omega_i$ be finite sets and let
$L_i\subseteq\R^{\Omega_i}$ be real linear spaces containing the constant
function.  On $\Omega=\prod_i\Omega_i$, define the additive degree-one space
\[
 L=\left\{c+\sum_i f_i(x_i):c\in\R,\ f_i\in L_i\right\}.
\]

For each factor, let $\mathcal Q_i$ be a set of nonconstant Boolean atoms,
with one representative from each complementary pair.  For a nonempty
finite $\mathcal F\subseteq\mathcal Q_i$, define its centered rank
\[
 \rho_i(\mathcal F)
 =\dim\operatorname{span}(\{1\}\cup\mathcal F)-1.
\]

\begin{theorem}[Atomic product transfer]\label{thm:atomic-transfer}
Suppose the following hold for fixed $C\ge1$ and $\alpha\ge1$:
\begin{enumerate}
\item every nonconstant Boolean function in $L_i$ is an atom in
      $\mathcal Q_i$ or its complement;
\item every nonempty finite $\mathcal F\subseteq\mathcal Q_i$ satisfies
      $|\mathcal F|\le C\rho_i(\mathcal F)^\alpha$.
\end{enumerate}
Then every Boolean matrix $M$ with column set $\Omega$ and rows in $L$
satisfies
\[
 \Dcc(M)\le
 \left\lceil\log_2\bigl(2C(\rank_{\R}M+1)^\alpha+2\bigr)\right\rceil.
\]
In particular, $\Dcc(M)\le\alpha\log_2(\rank M+1)+O(\log C+1)$.
\end{theorem}

\begin{proof}
Write a Boolean row as $f=c+\sum_i f_i(x_i)$.  Fixing all factors except
$i$ shows that the range of every nonconstant $f_i$ has exactly two values.
If two factors were nonconstant, the sum of their two two-point ranges would
contain at least three values, contradicting Booleanity.  Thus every
nonconstant row is a lifted atom from one factor or its complement.

Let $\mathcal F_i$ be the positive atoms used by the matrix rows and put
$K=\sum_i|\mathcal F_i|$.  Replacing complements by positive atoms enlarges
the row span by at most the constant function.  Hence all lifted atoms lie in
a space whose quotient by constants has dimension at most
$r=\rank_{\R}M$.

The centered one-factor spaces form a direct sum on the product: if
$\sum_i g_i(x_i)$ is constant, varying one coordinate shows that every
$g_i$ is constant.  Consequently
\[
 \sum_i\rho_i(\mathcal F_i)\le r.
\]
The atom-growth hypothesis and convexity give
\[
 K\le C\sum_i\rho_i(\mathcal F_i)^\alpha
 \le C r^\alpha.
\]
Alice sends one of two constant messages, or the used atom index and a
complement bit.  There are at most $2K+2$ messages, proving the theorem.
\end{proof}

The polynomial atom-growth hypothesis has a natural opposite obstruction.

\begin{proposition}[Selector universality]\label{prop:selector}
Suppose one factor contains Boolean degree-one functions
$q_1,\ldots,q_N$ satisfying
\[
 q_iq_j=0\ (i\ne j),
 \qquad
 \sum_{j=1}^Nq_j=1,
\]
and every selector value $j$ occurs on the domain.  Then row-degree-one
Boolean matrices on that factor contain every Boolean matrix with $N$
column types.
\end{proposition}

\begin{proof}
Choose one domain point $x_j$ with $q_j(x_j)=1$.  For an arbitrary Boolean
row vector $v\in\{0,1\}^N$, the degree-one function
\[
 f_v=\sum_{j:v_j=1}q_j
\]
is Boolean and satisfies $f_v(x_j)=v_j$.  Applying this construction to every
row embeds the prescribed matrix.  Extra domain points only duplicate one of
the selector column types.
\end{proof}

Thus polynomial centered-rank atom growth is a sufficient tractable regime,
while a large selector partition recovers unrestricted Boolean matrices.

\section{Low-degree cube and slice matrices}
\label{sec:low-degree-slice}

We next give a second transfer mechanism which is not limited to degree one.
For a Boolean function on the cube, degree means the degree of its unique
multilinear real polynomial.  For a function on $\binom{[n]}k$, slice degree
is the minimum degree of a real polynomial agreeing with it on the layer.

We use three published ingredients.  Filmus--Ihringer prove that there is a
universal $C_0>1$ such that, when
$\min(k,n-k)\ge C_0^s$, every Boolean slice function of degree at most $s$
is the restriction of a Boolean cube function of the same degree
\cite{FilmusIhringerJunta}.  Chiarelli--Hatami--Saks prove that every Boolean
cube function of degree at most $s$ depends on at most
$6.614\cdot2^s$ coordinates \cite{ChiarelliHatamiSaks}.  Finally, Midrijanis
proves that such a function has deterministic decision-tree depth at most
$2s^3$ \cite{Midrijanis}.  Put
\[
 J_s=7\cdot2^s,
 \qquad q_s=2s^3,
\]
and enlarge $C_0$, if necessary, so that $C_0^s\ge J_s$ for every $s\ge1$.

The new point is that the coordinate names used by the query trees can be
compressed in terms of the actual matrix rank.

\begin{lemma}[Rank pooling on the cube]
\label{lem:cube-rank-pooling}
Let $M$ be a Boolean matrix of rank $r\ge1$ whose columns are indexed by
$\{0,1\}^n$.  If every row depends on at most $J$ coordinates, then all rows
depend on one common set of at most $Jr$ coordinates.
\end{lemma}

\begin{proof}
Choose $r$ rows $f_1,\ldots,f_r$ forming a basis of the row space, choose a
relevant-coordinate set $S_i$ of size at most $J$ for each one, and put
$S=\bigcup_iS_i$.  If two columns agree on $S$, every basis row takes the same
value on them.  Every other row is a real linear combination of the basis
rows and therefore also takes the same value.  Thus every row depends on
$S$, while $|S|\le Jr$.
\end{proof}

\begin{theorem}[Low-degree cube Log-Rank theorem]
\label{thm:low-degree-cube}
Let $M$ be a Boolean matrix whose columns are indexed by $\{0,1\}^n$, and
suppose every row has degree at most $s\ge1$.  If
$r=\rank_{\R}M$, then
\[
 \Dcc(M)\le
 2s^3\left(
 \left\lceil\log_2\bigl(7\cdot2^s(r+1)\bigr)\right\rceil+1
 \right).
\]
\end{theorem}

\begin{proof}
If $r=0$, the matrix is identically zero.  Otherwise apply
Lemma~\ref{lem:cube-rank-pooling} with $J=J_s$.  For each row, Alice simulates
a depth-$q_s$ decision tree.  At every internal node she sends the index,
inside the common coordinate pool, of the next queried coordinate; Bob
returns its bit.  Each round costs at most
$\lceil\log_2(J_sr)\rceil+1$ bits.  Alice knows the reached leaf and outputs
its label.  Replacing $r$ by $r+1$ gives the displayed uniform bound.
\end{proof}

The only extra issue on a slice is that junta supports need not be unique,
because the total weight is fixed.  The following synchronization lemma
handles exactly this point.

\begin{lemma}[Slice support intersection]
\label{lem:slice-support-intersection}
Let $f:\binom{[n]}k\to\R$ be both an $S$-junta and a $T$-junta.  If
\[
 [n]\setminus(S\cup T)\ne\varnothing,
\]
then $f$ is an $(S\cap T)$-junta.
\end{lemma}

\begin{proof}
Put $U=S\cap T$ and choose $j\notin S\cup T$.  We first show that $f$ is
invariant under every transposition of two coordinates outside $U$.  A
transposition of two points outside $S$ preserves the $S$-trace and one of
two points outside $T$ preserves the $T$-trace.  The only remaining case,
after exchanging the two names, has
$a\in S\setminus T$ and $b\in T\setminus S$.  The identity
\[
 (a\ b)=(a\ j)(b\ j)(a\ j)
\]
expresses it as two transpositions outside $T$ and one outside $S$, so it
also preserves $f$.

If two $k$-sets have the same trace on $U$, their two symmetric-difference
halves have the same size and lie outside $U$.  Pair the two halves and swap
the paired coordinates.  The resulting sequence of transpositions outside
$U$ carries one set to the other, and every step preserves $f$.  Hence $f$
depends only on its $U$-trace.
\end{proof}

This lemma combines with the sharp fixed-degree junta threshold on the slice.

\begin{theorem}[Optimal-threshold fixed-degree slice theorem]
\label{thm:fixed-degree-slice}
For every $s\ge1$ there is an integer $m_s\ge1$ such that the following
holds.  If
\[
 \min(k,n-k)\ge2s
\]
and every row of a Boolean matrix $M$ on $\binom{[n]}k$ has slice degree at
most $s$, then, writing $r=\rank_{\R}M$,
\[
 \Dcc(M)\le
 \left\lceil\log_2\left(
 (m_s+1)\bigl(m_s(r+1)+1\bigr)^{m_s}2^{2^{m_s}}
 \right)\right\rceil.
\]
In particular, $\Dcc(M)=O_s(\log(r+1))$.
\end{theorem}

\begin{proof}
Filmus proves that every Boolean degree-$s$ function on such a slice is an
$m_s$-junta, for a constant $m_s$ depending only on $s$
\cite{FilmusJuntaThreshold}.  His theorem is stated for $k\le n/2$; the
other half follows by complementing every coordinate, which preserves
degree.

Choose $r$ basis rows and let $S$ be the union of junta supports of size at
most $m_s$ for those rows.  Then $|S|\le m_sr$, and every matrix row is an
$S$-junta.  An arbitrary row also has some junta support $T$ with
$|T|\le m_s$.  If $n>m_s(r+1)$, then $S\cup T\ne[n]$, and
Lemma~\ref{lem:slice-support-intersection} replaces $T$ by $S\cap T$.
If $n\le m_s(r+1)$, use all of $[n]$.  Thus in both cases every row has a
junta support of size at most $m_s$ inside one common universe $U$ with
\[
 |U|\le m_s(r+1).
\]

There are at most
\[
 \sum_{j=0}^{m_s}\binom{|U|}{j}2^{2^j}
 \le (m_s+1)(|U|+1)^{m_s}2^{2^{m_s}}
\]
distinct rows: choose a support and then a Boolean truth table on it.  Alice
sends the index of her row among these possibilities, and Bob evaluates it
on his column.  The displayed count proves the theorem.
\end{proof}

The preceding theorem uses a non-explicit junta-size constant.  A polynomial
layer-depth assumption permits a completely explicit bound.  The extra
ingredient is a symmetry lemma for decision trees.

\begin{lemma}[Symmetric-tree canonicalization]
\label{lem:symmetric-tree}
Let $f:\binom{[n]}k\to\{0,1\}$ be an $S$-junta and suppose $f$ has a
coordinate decision tree of depth at most $q$.  There is a set
$R\subseteq[n]\setminus S$ with $|R|\le q$ and a depth-$q$ decision tree for
$f$ all of whose queried coordinates lie in $S\cup R$.
\end{lemma}

\begin{proof}
Choose any $R\subseteq[n]\setminus S$ of size
$\min\{q,|[n]\setminus S|\}$.  We prove by induction on $q$ that every
depth-$q$ tree can be replaced by one using only $S\cup R$.

There is nothing to prove for a leaf.  Otherwise let $i$ be the root query.
If $i\notin S$ and $i\notin R$, choose $j\in R$ and conjugate the entire tree
by the transposition $(i\ j)$.  This still computes $f$, since the
transposition fixes $S$ and $f$ is an $S$-junta, and its root query is now
$j$.  Thus the root lies in $S\cup R$.

If the root is a point of $R$, remove it from $R$ for the recursive calls and
add it to the fixed support; each restricted function is invariant under all
permutations outside $S$ together with that root.  If the root lies in $S$,
use the same fixed subset of at most $q-1$ points of $R$ in both recursive
calls.  Conjugating subtrees in the same way proves the induction step.
\end{proof}

We can now state the general communication principle.  A coordinate domain
$\Omega\subseteq[Q]^n$ is called symmetric if it is invariant under every
permutation of its $n$ coordinates.

\begin{theorem}[Rank-pooled bounded-answer query transfer]
\label{thm:rank-pooled-query}
Let $M$ be a Boolean matrix of real rank $r$ whose column domain is either
\begin{enumerate}
\item a symmetric subdomain of $[Q]^n$, or
\item a complete product $\prod_{i=1}^n\Sigma_i$ with
      $2\le|\Sigma_i|\le Q$.
\end{enumerate}
Suppose every row function has a coordinate decision tree of depth at most
$q$.  Put $L=\lceil\log_2Q\rceil$.  Then
\[
 \Dcc(M)\le
 q\left(\lceil\log_2(r+1)\rceil+(q+1)L\right).
\]
Consequently, if both $q$ and $\log Q$ are polynomial in
$\log(r+1)$, then the Log-Rank Conjecture holds for $M$.
\end{theorem}

\begin{proof}
Assume $r\ge1$ and choose $r$ basis rows together with depth-$q$ coordinate
trees.  A $Q$-ary depth-$q$ tree has fewer than $Q^q$ internal nodes.  If $S$
is the union of every coordinate queried by the basis trees, then
\[
 |S|\le rQ^q.
\]
Every matrix row is a real linear combination of the basis rows and hence is
an $S$-junta.

In the symmetric case, Lemma~\ref{lem:symmetric-tree}, whose proof is
unchanged for a multiway tree, puts every row tree inside one common universe
$U=S\cup R$ with $|R|\le q$.  Thus
\[
 |U|\le rQ^q+q\le(r+1)Q^q.
\]
In the complete-product case, every query outside $S$ can instead be pruned:
all its answer restrictions compute the same residual function, and any one
child is valid on every assignment of the remaining coordinates.  Hence one
may take $U=S$.

Alice simulates the resulting row tree.  Each round she names a coordinate
of $U$, using at most
$\lceil\log_2(r+1)\rceil+qL$ bits, and Bob returns its value using $L$ bits.
There are at most $q$ rounds.
\end{proof}

Dafni, Filmus, Lifshitz, Lindzey, and Vinyals prove quantitative polynomial
relations between degree and coordinate decision-tree depth on multislices
\cite{DafniFilmusLifshitzLindzeyVinyals}.  Their parameters give the following
explicit slice consequence.  If
\[
 \min(k,n-k)\ge144s^4,
\]
then every Boolean slice function of degree at most $s$ has decision-tree
depth at most
\begin{equation}\label{eq:slice-query-depth}
 q_s=5184s^8.
\end{equation}
Indeed, in their notation a slice has $\Delta=1$, conflict bound
$M=\lfloor\min(k,n-k)/2\rfloor$, and block-sensitivity ratio
$\beta\ge1/2$.  Their Theorem~3.1 gives
$\operatorname{bs}(f)\le6s^2$ and, once
$s\le(\beta M/36)^{1/4}$,
\[
 D_{\rm query}(f)
 \le\beta^{-2}\Delta\operatorname{bs}(f)^4
 \le4(6s^2)^4=5184s^8.
\]

\begin{theorem}[Polynomial-depth low-degree slice Log-Rank theorem]
\label{thm:polynomial-depth-slice}
Let $M$ be a Boolean matrix whose columns are indexed by $\binom{[n]}k$ and
whose rows have slice degree at most $s\ge1$.  If
\[
 \min(k,n-k)\ge144s^4,
\]
put $r=\rank_{\R}M$ and $q_s=5184s^8$.  Then
\[
 \Dcc(M)\le
 q_s\bigl(\lceil\log_2(r+1)\rceil+q_s+1\bigr).
\]
In particular,
\[
 \Dcc(M)=O\bigl(s^8\log(r+1)+s^{16}\bigr).
\]
\end{theorem}

\begin{proof}
The Johnson slice is a symmetric binary coordinate domain.  Equation
\eqref{eq:slice-query-depth} supplies depth $q_s$ for every row.  Apply
Theorem~\ref{thm:rank-pooled-query} with $Q=2$ and $L=1$.
\end{proof}

\begin{theorem}[Very-deep slice refinement]
\label{thm:low-degree-slice}
There is a universal constant $C>1$ such that the following holds.  Let $M$
be a Boolean matrix whose columns are indexed by $\binom{[n]}k$, and suppose
every row has slice degree at most $s\ge1$.  If
\[
 \min(k,n-k)\ge C^s
\]
and $r=\rank_{\R}M$, then
\[
 \Dcc(M)\le
 2s^3\left(
 \left\lceil\log_2\bigl(7\cdot2^s(r+1)\bigr)\right\rceil+1
 \right).
\]
Consequently
\[
 \rect(M)\ge
 2^{-2s^3(\lceil\log_2(7\cdot2^s(r+1))\rceil+1)}.
\]
\end{theorem}

\begin{proof}
The rank-zero case is immediate, so assume $r\ge1$.  Choose a row-space
basis $f_1,\ldots,f_r$.  Extend each $f_i$ to a Boolean degree-$s$ cube
function and let $S_i$ be a relevant-coordinate set of size at most $J_s$.
Put $S=\bigcup_iS_i$, so $|S|\le J_sr$.  On the slice, every matrix row is a
linear combination of the $f_i$ and hence is an $S$-junta.

For an arbitrary row $f$, choose its own Boolean degree-$s$ cube extension
$g$ and a relevant set $T$ with $|T|\le J_s$.  If
$n>J_s(r+1)$, then $[n]\setminus(S\cup T)$ is nonempty.
Lemma~\ref{lem:slice-support-intersection} shows that the restriction of $g$
to the slice is an $(S\cap T)$-junta.  Every Boolean pattern on $T$ occurs on
the slice because $|T|\le J_s\le\min(k,n-k)$.  It follows that the total cube
function $g$ itself is independent of $T\setminus S$.  Thus its decision tree
queries only coordinates in $S$.

If instead $n\le J_s(r+1)$, use all of $[n]$ as the coordinate pool.  In both
cases every row has a depth-$q_s$ decision tree whose queried coordinates lie
in one common universe $U$ satisfying
\[
 |U|\le J_s(r+1).
\]
The protocol from Theorem~\ref{thm:low-degree-cube} now proves the
communication bound.  A deterministic $c$-bit protocol has at most $2^c$
transcript rectangles, so one monochromatic rectangle has relative density
at least $2^{-c}$.
\end{proof}

Theorem~\ref{thm:fixed-degree-slice} gives
$\Dcc(M)=O_s(\log(r+1))$ already at the optimal junta threshold.  The
polynomial-depth Theorem~\ref{thm:polynomial-depth-slice} gives a fully
explicit polylogarithmic bound whenever $s=\operatorname{polylog}(r+1)$,
under the polynomial layer-depth condition $144s^4$.  In the still deeper
$C^s$ regime, Theorem~\ref{thm:low-degree-slice} improves the dependence on
$s$ from degrees eight and sixteen to degrees three and four.  All layer-depth
conditions concern only the ambient slice and are independent of matrix rank.

The opposite boundary is universal.

\begin{proposition}[Low-degree boundary universality]
\label{prop:low-degree-boundary}
If $k\le s$, every Boolean function on $\binom{[n]}k$ has slice degree at
most $s$.  The same holds if $n-k\le s$.  In particular, the case $k=1$
contains every Boolean matrix with $n$ columns.
\end{proposition}

\begin{proof}
For $\mathcal F\subseteq\binom{[n]}k$, its indicator on the layer is
\[
 \1_{\mathcal F}(x)=
 \sum_{A\in\mathcal F}\prod_{i\in A}x_i.
\]
Exactly one degree-$k$ monomial is one at each layer point.  This proves the
first assertion; the second follows by replacing every coordinate by its
complement.  When $k=1$, the columns are the coordinate unit vectors and an
arbitrary row is the indicator of an arbitrary subset of them.
\end{proof}

The threshold $2s$ in Theorem~\ref{thm:fixed-degree-slice} is sharp for a
uniform bounded-junta conclusion \cite{FilmusJuntaThreshold}.  The polynomial
$144s^4$ hypothesis supplies an explicit shallow decision tree, while the
larger $C^s$ hypothesis supplies a same-degree total-cube extension and the
sharper dependence on $s$.

\subsection{Bounded alphabets and multislices}

The rank-pooling mechanism extends beyond binary slices.  Let
\[
 \Omega_{\rm prod}=\prod_{i=1}^n[m_i],
 \qquad 2\le m_i\le Q,
\]
and define degree using the coordinate-value indicators
$\1_{\{x_i=a\}}$.  A query asks for $x_i$ and has at most $Q$ answers.

\begin{theorem}[Low-degree bounded-alphabet product theorem]
\label{thm:low-degree-product-domain}
Let $M$ be a Boolean matrix with column domain $\Omega_{\rm prod}$ and rows
of degree at most $s\ge1$.  Put
\[
 r=\rank_{\R}M,
 \qquad q=1296s^8,
 \qquad L=\lceil\log_2Q\rceil.
\]
Then
\[
 \Dcc(M)\le
 q\left(\lceil\log_2(r+1)\rceil+(q+1)L\right).
\]
For fixed $Q$, this is
$O_Q(s^8\log(r+1)+s^{16})$.
\end{theorem}

\begin{proof}
For a complete product domain, the parameters in
\cite[Theorem~3.1]{DafniFilmusLifshitzLindzeyVinyals} are
$\Delta=1$, conflict bound $M_0=n$, and $\beta=1$.  Hence every Boolean
degree-$s$ row has coordinate decision-tree depth at most
\[
 \operatorname{bs}(f)^4\le(6s^2)^4=1296s^8=q.
\]
Apply Theorem~\ref{thm:rank-pooled-query} in its complete-product case.
\end{proof}

The bounded-alphabet condition is essential.  A single coordinate with a
large alphabet is a selector domain: its degree-one Boolean functions already
encode arbitrary Boolean matrices.

Now let $\boldsymbol\lambda=(\lambda_1,\ldots,\lambda_m)$ be positive
integers summing to $n$, with $m\le Q$, and consider the multislice
\[
 \Omega(\boldsymbol\lambda)
 =\{x\in[m]^n:|\{i:x_i=c\}|=\lambda_c\text{ for every }c\}.
\]
Write $\lambda_{\min}=\min_c\lambda_c$.

\begin{theorem}[Low-degree multislice Log-Rank theorem]
\label{thm:low-degree-multislice}
Let $M$ be a Boolean matrix with column domain
$\Omega(\boldsymbol\lambda)$ and rows of multislice degree at most $s\ge1$.
Assume
\[
 \lambda_{\min}\ge216s^4.
\]
Put
\[
 r=\rank_{\R}M,
 \qquad q=11664s^8,
 \qquad L=\lceil\log_2Q\rceil.
\]
Then
\[
 \Dcc(M)\le
 q\left(\lceil\log_2(r+1)\rceil+(q+1)L\right).
\]
For fixed $Q$, this is
$O_Q(s^8\log(r+1)+s^{16})$.
\end{theorem}

\begin{proof}
For a multislice, the parameters in
\cite[Theorem~3.1 and Section~4.3]{DafniFilmusLifshitzLindzeyVinyals} are
\[
 \Delta=1,
 \qquad M_0=\lfloor\lambda_{\min}/2\rfloor,
 \qquad \beta\ge1/3.
\]
The layer hypothesis implies
$s\le(\beta M_0/36)^{1/4}$.  Their strengthened bound therefore gives
\[
 D_{\rm query}(f)
 \le\beta^{-2}\operatorname{bs}(f)^4
 \le9(6s^2)^4=11664s^8=q.
\]
The multislice is symmetric under coordinate permutations.  Apply
Theorem~\ref{thm:rank-pooled-query}.
\end{proof}

When a color has multiplicity one, degree-one selector universality already
recovers arbitrary Boolean matrices.  Thus some lower bound on
$\lambda_{\min}$ is structurally unavoidable, although the polynomial
$216s^4$ threshold is not claimed to be sharp.

\subsection{Unisolvent sparse-column transfer}

The preceding theorems remain valid after a sparse column restriction that
preserves the ambient function space.  Let $V\subseteq\R^\Omega$ be a
finite-dimensional space and $Y\subseteq\Omega$.  Call $Y$
\emph{$V$-unisolvent} if the restriction map
\[
 \operatorname{res}_Y:V\longrightarrow\R^Y
\]
is injective.

\begin{theorem}[Unisolvent sparse-column transfer]
\label{thm:unisolvent-transfer}
Assume $\Omega$ satisfies one of the two domain hypotheses in
Theorem~\ref{thm:rank-pooled-query}, every Boolean function in $V$ has
coordinate-query depth at most $q$, and $Y$ is $V$-unisolvent.  Then every
Boolean matrix $M$ with columns $Y$ whose rows extend to Boolean functions in
$V$ satisfies
\[
 \Dcc(M)\le
 q\left(\lceil\log_2(\rank_{\R}M+1)\rceil+
 (q+1)\lceil\log_2Q\rceil\right).
\]
\end{theorem}

\begin{proof}
Let $\widetilde M$ be the matrix with the same row functions and all ambient
columns $\Omega$.  A row relation of $\widetilde M$ remains one after
restriction.  Conversely, a restricted relation is an element of $V$
vanishing on $Y$, and unisolvence forces it to vanish on $\Omega$.  Hence
\[
 \rank_{\R}\widetilde M=\rank_{\R}M.
\]
Apply Theorem~\ref{thm:rank-pooled-query} to $\widetilde M$ and restrict Bob's
inputs to $Y$.
\end{proof}

The assumption that each sparse row already has a Boolean ambient extension
can itself be certified by unisolvence at twice the degree.

\begin{theorem}[Booleanity-certifying sparse transfer]
\label{thm:booleanity-certifying-transfer}
Let $V_{\le s}\subseteq V_{\le2s}\subseteq\R^\Omega$ be a degree filtration
such that
\[
 f\in V_{\le s}\quad\Longrightarrow\quad f(f-1)\in V_{\le2s}.
\]
If $Y$ is $V_{\le2s}$-unisolvent and $f\in V_{\le s}$ takes values in
$\{0,1\}$ on $Y$, then $f$ is Boolean on all of $\Omega$.  Consequently
Theorem~\ref{thm:unisolvent-transfer} applies to Boolean matrices on $Y$
whose rows merely extend to real functions in $V_{\le s}$.
\end{theorem}

\begin{proof}
The function $h=f(f-1)$ belongs to $V_{\le2s}$ and vanishes on $Y$.
Unisolvence gives $h=0$ on $\Omega$, so $f(x)\in\{0,1\}$ everywhere.
Moreover $V_{\le s}\subseteq V_{\le2s}$, so every restricted row relation
lifts to the ambient domain as in Theorem~\ref{thm:unisolvent-transfer}.
\end{proof}

We next show that random sparse subsets are unisolvent on every transitive
invariant function space.  Let a finite group act transitively on $\Omega$,
let $V\subseteq\R^\Omega$ be invariant, and put $d=\dim V$.  Choose an
orthonormal basis $\phi_1,\ldots,\phi_d$ for the uniform inner product and
write
\[
 z_x=(\phi_1(x),\ldots,\phi_d(x))^\top.
\]
The diagonal of the reproducing kernel is invariant and hence constant.
Taking its average gives
\begin{equation}\label{eq:constant-leverage}
 \|z_x\|_2^2=d,
 \qquad
 \mathbb E_{x\in\Omega}z_xz_x^\top=I_d.
\end{equation}

\begin{theorem}[Random unisolvent sampling]
\label{thm:random-unisolvent}
Let $Y$ be a uniformly random $N$-element subset of $\Omega$, sampled
without replacement.  Then
\[
 \Pr[Y\text{ is not }V\text{-unisolvent}]
 \le d\exp\left(-\frac{N}{8d}\right).
\]
Consequently $N\ge8d\log(d/\delta)$ gives unisolvence with probability at
least $1-\delta$.
\end{theorem}

\begin{proof}
Restriction is injective exactly when
\[
 G_Y=\sum_{x\in Y}z_xz_x^\top
\]
is positive definite.  By \eqref{eq:constant-leverage}, the summands are
positive semidefinite, bounded by $dI_d$, and have mean $I_d$.  Matrix
Chernoff at lower deviation $1/2$ \cite{Tropp}, together with the
without-replacement Laplace-transform comparison of Gross--Nesme
\cite{GrossNesme}, yields
\[
 \Pr[\lambda_{\min}(G_Y)\le N/2]
 \le d\left(\sqrt{\frac2e}\right)^{N/d}
 \le d e^{-N/(8d)}.
\]
Outside this event $G_Y$ is positive definite.
\end{proof}

The logarithmic oversampling factor can be removed by standard projection-DPP
volume sampling.  Let $A\in\R^{|\Omega|\times d}$ have orthonormal columns
spanning the evaluation image of $V$ in the standard Euclidean inner product,
and put $P=AA^\top$.

\begin{theorem}[Minimal projection-DPP unisolvent sampling]
\label{thm:dpp-unisolvent}
For every $d$-element subset $Y\subseteq\Omega$, define
\[
 \Pr_{\rm DPP}[Y]=\det(P_Y).
\]
These probabilities sum to one, and every set in their support is
$V$-unisolvent.  Thus the process always selects exactly the
information-theoretic minimum $d=\dim V$ columns.
\end{theorem}

\begin{proof}
For $|Y|=d$,
\[
 \det(P_Y)=\det(A_YA_Y^\top)=\det(A_Y)^2.
\]
Cauchy--Binet gives
\[
 \sum_{|Y|=d}\det(P_Y)
 =\sum_{|Y|=d}\det(A_Y)^2
 =\det(A^\top A)=1.
\]
Positive probability is equivalent to $A_Y$ being nonsingular, which is
exactly unisolvence.  Conversely, injectivity of $V\to\R^Y$ requires
$|Y|\ge d$.
\end{proof}

When $V$ is invariant under a transitive group, $P$ and the resulting
projection DPP are invariant.  Its one-point marginals are uniform:
\[
 \Pr[x\in Y]=P_{xx}=d/|\Omega|.
\]

For the Johnson slice, let $V_{\le s}$ be the space of functions of slice
degree at most $s$.  The harmonic decomposition gives
\begin{equation}\label{eq:slice-degree-dimension}
 \dim V_{\le s}
 =\sum_{j=0}^s\left(\binom nj-\binom n{j-1}\right)
 =\binom ns
 \qquad(s\le\min(k,n-k))
\end{equation}
\cite{FilmusOrthogonal}, with $\binom n{-1}=0$.  This space is invariant
under the transitive $S_n$-action.

\begin{corollary}[Exponentially sparse random Johnson columns]
\label{cor:sparse-random-johnson}
Assume
\[
 \min(k,n-k)\ge144s^4,
 \qquad d_s=\binom ns,
\]
and choose a uniformly random $N$-element family
$Y\subseteq\binom{[n]}k$ with
\[
 N\ge8d_s\log(d_s/\delta).
\]
With probability at least $1-\delta$, simultaneously for every Boolean
matrix whose rows extend to degree-$s$ Boolean functions on the full slice,
\[
 \Dcc(M)\le
 q_s\bigl(\lceil\log_2(\rank_{\R}M+1)\rceil+q_s+1\bigr),
 \qquad q_s=5184s^8.
\]
\end{corollary}

\begin{proof}
Theorem~\ref{thm:random-unisolvent} and
\eqref{eq:slice-degree-dimension} make $Y$ $V_{\le s}$-unisolvent with the
stated probability.  Equation~\eqref{eq:slice-query-depth} supplies query
depth $q_s$, so Theorem~\ref{thm:unisolvent-transfer} applies with $Q=2$.
The unisolvence event depends only on $Y$, which gives the simultaneous
statement over all later row families.
\end{proof}

\begin{corollary}[Minimal projection-DPP Johnson columns]
\label{cor:dpp-johnson}
Under the hypotheses of Corollary~\ref{cor:sparse-random-johnson}, sample
$Y$ from the rank-$\binom ns$ projection DPP of $V_{\le s}$.  Then
\[
 |Y|=\binom ns
\]
and, with probability one, the communication conclusion of that corollary
holds simultaneously for every Boolean degree-$s$ row family restricted to
$Y$.  No smaller column set can be unisolvent for the entire space
$V_{\le s}$.
\end{corollary}

\begin{proof}
Apply Theorem~\ref{thm:dpp-unisolvent}, then
Theorem~\ref{thm:unisolvent-transfer}.
\end{proof}

The twice-degree certificate gives a particularly sharp adjacent-two-value
theorem with only quadratically many columns.

\begin{theorem}[Explicit quadratic unisolvent cross]
\label{thm:explicit-cross}
Assume $2\le k\le n-2$.  Fix a $k$-set $I$, put $O=[n]\setminus I$, and
choose fixed pairs $I_0\in\binom I2$ and $O_0\in\binom O2$.  Define
\begin{align*}
 \cY_\times={}&\{I\}
 \cup\{I\setminus\{i\}\cup\{a\}:i\in I,\ a\in O\}\\
 &{}\cup\{I\setminus\{i,j\}\cup O_0:\{i,j\}\in\tbinom I2\}\\
 &{}\cup\{I\setminus I_0\cup\{a,b\}:\{a,b\}\in\tbinom O2\}.
\end{align*}
Then
\[
 |\cY_\times|=\binom n2,
\]
and $\cY_\times$ is unisolvent for the degree-at-most-two slice space.
\end{theorem}

\begin{proof}
The two double-swap families overlap only at $I\setminus I_0\cup O_0$, so
\[
 |\cY_\times|
 =1+k(n-k)+\binom k2+\binom{n-k}2-1
 =\binom n2.
\]
On the weight-$k$ slice, pair monomials span all degree-at-most-two functions:
\[
 x_i=\frac1{k-1}\sum_{j\ne i}x_ix_j,
 \qquad
 1=\frac1{\binom k2}\sum_{i<j}x_ix_j.
\]
It suffices to show that a pair polynomial
\[
 F(B)=\sum_{\{u,v\}\subseteq B}c_{uv}
\]
vanishing on $\cY_\times$ has all coefficients zero.

Write
\[
 S=\sum_{\{i,j\}\in\binom I2}c_{ij},
 \qquad X_i=\sum_{j\in I\setminus\{i\}}c_{ij},
 \qquad Y_a=\sum_{i\in I}c_{ia}.
\]
Evaluation at $I$ gives $S=0$ and hence $\sum_iX_i=0$.  Evaluation at
$I\setminus\{i\}\cup\{a\}$ gives
\[
 c_{ia}=Y_a-X_i.
\]
Summing over $i$ yields $(k-1)Y_a=0$, so
\begin{equation}\label{eq:cross-coefficients}
 Y_a=0,
 \qquad c_{ia}=-X_i.
\end{equation}

Write $O_0=\{a_0,a_1\}$ and $z_0=c_{a_0a_1}$.  Evaluation at
$I\setminus\{i,j\}\cup O_0$, using \eqref{eq:cross-coefficients}, gives
\begin{equation}\label{eq:inside-coefficients}
 c_{ij}=-X_i-X_j-z_0.
\end{equation}
Summing \eqref{eq:inside-coefficients} over $j\ne i$ and using
$\sum_jX_j=0$ gives $X_i=-z_0$ for every $i$.  Summing in $i$ gives
$z_0=0$, so all inside and cross coefficients vanish.  Finally, evaluation
at $I\setminus I_0\cup\{a,b\}$ gives $c_{ab}=0$ for every outside pair.
Thus $F=0$.
\end{proof}

\begin{theorem}[Sparse adjacent-two-value Log-Rank from unisolvence]
\label{thm:quadratic-column-hls}
Assume $2\le k\le n-2$, and let $Y\subseteq\binom{[n]}k$ be unisolvent for
the degree-at-most-two slice space.  Then the following holds simultaneously
for every $\cA\subseteq2^{[n]}$ and integer $a$ satisfying
\[
 |A\cap B|\in\{a,a+1\}
 \qquad(A\in\cA,\ B\in Y).
\]
For $M_{A,B}=|A\cap B|-a$ and $r=\rank_{\R}M$,
\[
 \Dcc(M)\le\lceil\log_2(r+1)\rceil+1,
 \qquad
 \rect(M)\ge\frac1{4(r+1)}.
\]
\end{theorem}

\begin{proof}
For fixed $A$, the affine function $f_A(B)=|A\cap B|-a$ has slice degree at
most one and is Boolean on $Y$.  The selected set is unisolvent through
degree two by hypothesis, so Theorem~\ref{thm:booleanity-certifying-transfer} makes $f_A$
Boolean on the complete slice.  Every Boolean degree-one function on an
interior Johnson slice is constant, a coordinate query, or its complement
\cite{FilmusIhringer}.

Choose $r$ basis rows.  Their depth-one trees query at most $r$ coordinates,
and every other row is a junta over this pooled set.  The proof of
Lemma~\ref{lem:symmetric-tree} uses at most one additional representative.
Alice sends its index among at most $r+1$ possibilities and Bob returns one
bit, proving the communication bound.  Finally,
\[
 2^{-\lceil\log_2(r+1)\rceil-1}\ge\frac1{4(r+1)},
\]
so a largest transcript rectangle gives the stated density.
\end{proof}

\begin{corollary}[Quadratic-column adjacent-two-value Log-Rank]
\label{cor:quadratic-column-hls}
For every $2\le k\le n-2$, there exists a fixed family
$Y\subseteq\binom{[n]}k$ with
\[
 |Y|=\binom n2
\]
for which all conclusions of Theorem~\ref{thm:quadratic-column-hls} hold.
Moreover, a sample from the degree-two projection DPP has this property with
probability one.
\end{corollary}

\begin{proof}
Take $Y=\cY_\times$ from Theorem~\ref{thm:explicit-cross}.  The final DPP
assertion follows from Theorem~\ref{thm:dpp-unisolvent} and
\eqref{eq:slice-degree-dimension}.
\end{proof}

On balanced layers the column count can be reduced further, from quadratic
to linear, by separating Booleanity detection from rank preservation.  Put
\begin{equation}\label{eq:four-fiber-gap}
 \gamma_{n,k}
 =\frac{\binom{n-4}{k-2}}{\binom nk}
 =\frac{k(k-1)(n-k)(n-k-1)}{n(n-1)(n-2)(n-3)}.
\end{equation}

\begin{lemma}[Four-coordinate two-value gap]
\label{lem:four-coordinate-gap}
For every $A\subseteq[n]$ with $2\le|A|\le n-2$ and every integer $a$,
\[
 \Pr_{B\in\binom{[n]}k}[|A\cap B|\in\{a,a+1\}]
 \le1-\gamma_{n,k}.
\]
\end{lemma}

\begin{proof}
Choose $u,v\in A$ and $p,q\notin A$, all distinct, and partition the
$k$-sets by their trace outside $D=\{u,v,p,q\}$.  A fiber whose outside
trace has size $k-2$ contains the six two-subsets of $D$.  Their local
contributions to $|A\cap B|$ are $0,1,2$ with multiplicities $1,4,1$.
Any two consecutive integers omit at least one extreme.  There are
$\binom{n-4}{k-2}$ such fibers, proving the claim.
\end{proof}

The fourth-power parameter in Lemma~\ref{lem:four-coordinate-gap} is an
artifact of activating one fixed four-coordinate fiber.  We next prove a
quadratic bound.  Put
\begin{equation}\label{eq:quadratic-gap}
 p_{n,k}=\frac{\min\{k,n-k\}}n,
 \qquad
 \eta_{n,k}=\frac{p_{n,k}^2}{16},
 \qquad
 \xi_{n,k}=\frac25p_{n,k}^2,
 \qquad
 \theta_{n,k}=\max\{\gamma_{n,k},\xi_{n,k}\}.
\end{equation}

\begin{lemma}[Log-concave two-point escape]
\label{lem:log-concave-two-point}
Let $Z$ be integer-valued with a finite, interval-supported, log-concave
probability mass function.  If $\Var(Z)\ge1/3$, then for every integer $a$,
\[
 \Pr[Z\notin\{a,a+1\}]\ge\frac1{64}.
\]
\end{lemma}

\begin{proof}
Write $\varepsilon=\Pr[Z\notin\{a,a+1\}]$, $u=\Pr[Z=a]$, and
$v=\Pr[Z=a+1]$, and suppose that $\varepsilon<1/64$.  If
$M=\max\{u,v\}$ and $r=\sqrt{\varepsilon/M}$, then
\[
 M>63/128,
 \qquad r^2<2/63<1/25.
\]
Log-concavity and the fact that each first tail atom is at most
$\varepsilon$ give
\[
 \frac{\Pr[Z=a+2]}v\le r,
 \qquad
 \frac{\Pr[Z=a-1]}u\le r,
\]
with the zero cases interpreted using interval support.  Indeed, when
$u>v>0$, the first inequality follows from
\[
 \frac{\Pr[Z=a+2]}v
 \le\min\left\{\frac vu,\frac\varepsilon v\right\}
 \le\sqrt{\frac\varepsilon u},
\]
and the other order and the left tail are symmetric.  Successive mass ratios
decrease away from a log-concave mode.  Put
$t_-=\Pr[Z=a-1]$ and $t_+=\Pr[Z=a+2]$.  Hence for $j\ge1$ the two tail atoms
at distance $j+1/2$ from $a+1/2$ are at most
$t_-r^{j-1}$ and $t_+r^{j-1}$, respectively, while
$t_-+t_+\le\varepsilon$.

Since $r<1/5$,
\[
 \sum_{j\ge1}(j+1/2)^2r^{j-1}
 \le\frac{1+r}{(1-r)^3}+\frac1{(1-r)^2}+\frac1{4(1-r)}
 \le\frac{135}{32}<\frac{17}{4}.
\]
Centering at $a+1/2$ therefore gives
\[
 \Var(Z)
 \le\mathbb E(Z-a-1/2)^2
 <\frac14+\frac{17}{4}\varepsilon
 <\frac{81}{256}<\frac13,
\]
a contradiction.
\end{proof}

\begin{theorem}[Quadratic hypergeometric two-point gap]
\label{thm:quadratic-hypergeometric-gap}
For every $A\subseteq[n]$ with $2\le|A|\le n-2$ and every integer $a$,
\[
 \Pr_{B\in\binom{[n]}k}[|A\cap B|\notin\{a,a+1\}]
 \ge\eta_{n,k}.
\]
\end{theorem}

\begin{proof}
Complementing $B$ replaces $k$ by $n-k$ and $|A\cap B|$ by
$|A|-|A\cap B|$.  Complementing $A$ replaces $|A|$ by $n-|A|$ and the
intersection by $k-|A\cap B|$.  Thus assume
$2\le k,m:=|A|\le n/2$.  Set
\[
 p=k/n,
 \qquad q=m/n,
 \qquad \mu=\mathbb E|A\cap B|=pqn.
\]
The hypergeometric mass ratios
\[
 \frac{\Pr[X=j+1]}{\Pr[X=j]}
 =\frac{(m-j)(k-j)}{(j+1)(n-m-k+j+1)}
\]
decrease in $j$, so the law is interval-supported and log-concave.

Suppose first that $\mu\ge1$.  Since
\[
 p+q\le\frac12+2pq
\]
for $p,q\le1/2$, the exact hypergeometric variance satisfies
\[
 \Var(X)
 =\mu(1-p)(1-q)\frac n{n-1}
 \ge\frac{\mu(n-2\mu)}{2(n-1)}
 \ge\frac13.
\]
The last quadratic is concave on $[1,n/4]$, and both endpoint values are at
least $1/3$.  Lemma~\ref{lem:log-concave-two-point} gives escape at least
$1/64\ge p^2/16$.

Now suppose $\mu<1$.  For $a\le-2$ the conclusion is immediate.  For
$a=-1$, fixing one element of $A$ gives escape probability at least
$p\ge p^2/16$.  For $a=0$, fixing two elements gives
\[
 \Pr[X\ge2]\ge\frac{k(k-1)}{n(n-1)}\ge\frac{p^2}{2}.
\]
Finally let $a\ge1$.  Then $X=0$ is an escape value.  Integrality of
$km<n$ and $2\le k\le n/2$ gives, for $0\le i<k$,
\[
 \frac m{n-i}
 \le\frac{n-1}{k(n-k+1)}\le\frac12,
 \qquad
 \sum_{i=0}^{k-1}\frac m{n-i}<2.
\]
Here the bound by $1/2$ follows from
$k(n-k+1)-2(n-1)=(k-2)(n-k-1)\ge0$.
By concavity, $\log(1-x)\ge-x\log4$ on $[0,1/2]$.  Therefore
\[
 \Pr[X=0]
 =\prod_{i=0}^{k-1}\left(1-\frac m{n-i}\right)
 >\frac1{16}\ge\frac{p^2}{16}.
\]
This covers every adjacent pair.
\end{proof}

\begin{proposition}[Sharp \(2/5\) gap for row short-side at most three]
\label{prop:sharp-three-marked-gap}
Let
\[
 q=\min\{k,n-k\},
 \qquad
 s=\min\{|A|,n-|A|\}.
\]
If \(s\in\{2,3\}\), then for every integer \(a\),
\[
 \Pr_{B\in\binom{[n]}k}[|A\cap B|\notin\{a,a+1\}]
 \ge\frac25\left(\frac qn\right)^2.
\]
Equality holds if and only if
\[
 (n,k,|A|,a)=(6,3,3,1).
\]
\end{proposition}

\begin{proof}
Complementing \(A\) and \(B\) reduces to
\(2\le k\le n/2\) and \(m:=|A|\in\{2,3\}\).
For \(m=2\), the least mass left outside an adjacent pair is
\[
 \frac{k(k-1)}{n(n-1)}
 >\frac25\left(\frac kn\right)^2.
\]

Let \(m=3\), write \(P_j=\Pr[|A\cap B|=j]\), and put
\(D=n(n-1)(n-2)\).  Then
\[
\begin{aligned}
 P_0&=\frac{(n-k)(n-k-1)(n-k-2)}D,&
 P_1&=\frac{3k(n-k)(n-k-1)}D,\\
 P_2&=\frac{3k(k-1)(n-k)}D,&
 P_3&=\frac{k(k-1)(k-2)}D.
\end{aligned}
\]
Since \(P_0\ge P_3\) and \(P_1\ge P_2\), it is enough to bound
\(P_2+P_3\) and \(P_0+P_3\).

Write \(n=2k+t\), \(t\ge0\).  After clearing positive denominators, the first
bound is the strict positivity of
\[
\begin{aligned}
 H(k,t)={}&16k(k-1)(2k-1)\\
 &+(42k^2-54k+10)t+(13k-15)t^2,
\end{aligned}
\]
whose displayed coefficients are positive for \(k\ge2\).
The second bound is equivalent to \(F(k,t)\ge0\), where
\[
\begin{aligned}
 F(k,t)={}&12k^2(k-1)(k-3)
 +(32k^3-84k^2+40k)t\\
 &+(43k^2-60k+10)t^2+(25k-15)t^3+5t^4.
\end{aligned}
\]
For \(k\ge3\), every coefficient is nonnegative; the only possible zero is
\(k=3,t=0\).  For \(k=2\), the constraint \(m=3\le n/2\) gives \(t\ge2\),
and
\[
 F(2,t)=-48+62t^2+35t^3+5t^4>0.
\]
Thus equality occurs only at \(n=6,k=3\), where the escaping mass
\(P_0+P_3=1/10\) corresponds to the captured pair \(\{1,2\}\), hence \(a=1\).
\end{proof}




\begin{proposition}[Sharp \(2/5\) gap on the low-density side]
\label{prop:sharp-low-density-gap}
If
\[
 \frac{\min\{k,n-k\}}n\le\frac13,
\]
then for every nontrivial \(A\subseteq[n]\) and every integer \(a\),
\[
 \Pr_{B\in\binom{[n]}k}[|A\cap B|\notin\{a,a+1\}]
 >
 \frac25\left(\frac{\min\{k,n-k\}}n\right)^2.
\]
\end{proposition}

\begin{proof}
Complement \(B\) and write \(n=3k+t\), where \(k\ge2\) and \(t\ge0\).
After substituting the four-coordinate lower bound and clearing positive
denominators, the assertion is equivalent to \(L(k,t)>0\), where
\[
\begin{aligned}
L(k,t)={}&6k^4+18k^3-36k^2+12k
 +(26k^3-33k^2+3k)t\\
&+(17k^2-28k+5)t^2+(3k-5)t^3.
\end{aligned}
\]
Putting \(k=q+2\) gives
\[
\begin{aligned}
L(q+2,t)={}&6q^4+66q^3+216q^2+276q+120\\
&+(26q^3+123q^2+183q+82)t\\
&+(17q^2+40q+17)t^2+(3q+1)t^3,
\end{aligned}
\]
which is strictly positive for \(q,t\ge0\).
\end{proof}

Propositions~\ref{prop:sharp-three-marked-gap} and
\ref{prop:sharp-low-density-gap} reduce the unresolved sharp \(2/5\) theorem
to
\[
 \frac13<\frac{\min\{k,n-k\}}n\le\frac12,
 \qquad
 \min\{|A|,n-|A|\}\ge4.
\]



\begin{theorem}[Global sharp hypergeometric two-point gap]
\label{thm:global-sharp-hypergeometric-gap}
For every \(A\subseteq[n]\) with \(2\le|A|\le n-2\) and every integer \(a\),
\[
 \Pr_{B\in\binom{[n]}k}[|A\cap B|\notin\{a,a+1\}]
 \ge\frac25
 \left(\frac{\min\{k,n-k\}}n\right)^2.
\]
Equality holds if and only if
\[
 (n,k,|A|,a)=(6,3,3,1).
\]
\end{theorem}

\begin{proof}
Complement \(A\) and \(B\) so that \(2\le k\le n/2\) and
\(2\le m:=|A|\le n/2\).  The case \(m=2\) is strict by
Proposition~\ref{prop:sharp-three-marked-gap}.  For \(m\ge3\), fix three
points of \(A\) and three of its complement.  Partition the \(k\)-sets by
their trace outside these six coordinates.

For local trace sizes \(2,3,4\), the local intersection multiplicities are
\[
 (3,9,3),\qquad(1,9,9,1),\qquad(3,9,3).
\]
Every adjacent pair misses at least \(3,2,3\) local columns, respectively.
Summing the fibers gives escape probability at least
\[
 \zeta_{n,k}
 =
 \frac{3\binom{n-6}{k-2}
 +2\binom{n-6}{k-3}
 +3\binom{n-6}{k-4}}{\binom nk}.
\]
Here and below, out-of-range binomial coefficients are interpreted as zero.

Write \(n=2k+t\), \(t\ge0\).  Clearing the positive denominators in
\(\zeta_{n,k}\ge2(k/n)^2/5\) gives \(G(k,t)\ge0\).  For \(k=q+3\),
\[
\begin{aligned}
G(q+3,t)={}&16q^6+228q^5+1180q^4+2820q^3+3124q^2+1272q\\
&+(120q^5+1140q^4+4040q^3+6480q^2+4432q+876)t\\
&+(230q^4+1665q^3+4325q^2+4710q+1750)t^2\\
&+(195q^3+1065q^2+1925q+1140)t^3\\
&+(80q^2+300q+290)t^4+(13q+24)t^5.
\end{aligned}
\]
All coefficients are nonnegative, and the unique zero is \(q=t=0\).
For \(k=2\), one has \(t\ge2\); putting \(t=s+2\) gives
\[
G(2,s+2)
=11s^5+180s^4+1085s^3+3000s^2+3764s+1680>0.
\]
Thus equality requires \(n=6,k=3\).  Then the six-coordinate fiber is the
whole ground set, \(m=3\), and its multiplicities are \((1,9,9,1)\).
Exactly two of twenty columns escape only when the captured pair is
\(\{1,2\}\), giving \(a=1\).
\end{proof}


Combining Lemma~\ref{lem:four-coordinate-gap} and
Theorem~\ref{thm:global-sharp-hypergeometric-gap}, every nontrivial row escapes
every consecutive pair on at least a $\theta_{n,k}$-fraction of the layer.

We also need an explicit linear-size rank anchor.  Fix $i_0\in I$ and
$a_0\notin I$, and set
\begin{equation}\label{eq:linear-rank-anchor}
 \cY_0=\{I\}
 \cup\{I\setminus\{i_0\}\cup\{a\}:a\notin I\}
 \cup\{I\setminus\{i\}\cup\{a_0\}:i\in I\setminus\{i_0\}\}.
\end{equation}
Then $|\cY_0|=n$.

\begin{lemma}[Linear degree-one anchor]
\label{lem:linear-rank-anchor}
The family $\cY_0$ is unisolvent for the degree-at-most-one slice space.
\end{lemma}

\begin{proof}
If $f(B)=c+\sum_{j\in B}w_j$ vanishes on $\cY_0$, comparison with $I$ gives
$w_a=w_{i_0}$ for every $a\notin I$, and $w_i=w_{a_0}$ for every
$i\in I\setminus\{i_0\}$.  Thus all $w_j$ are equal, so $f$ is constant on
the slice; its value at $I$ is zero.
\end{proof}

\begin{theorem}[Linear-size sparse HLS on balanced layers]
\label{thm:linear-sparse-hls}
Let $\cY_1$ be a uniformly random $N$-element subset of
$\binom{[n]}k$, sampled without replacement, and put
$\cY=\cY_0\cup\cY_1$.  If
\[
 N\ge
 \frac{n\log2+\log(n+2)+\log(1/\delta)}{\theta_{n,k}},
\]
then, with probability at least $1-\delta$, simultaneously for every
$\cA\subseteq2^{[n]}$ and integer $a$ satisfying
\[
 |A\cap B|\in\{a,a+1\}
 \qquad(A\in\cA,\ B\in\cY),
\]
the matrix $M_{A,B}=|A\cap B|-a$ satisfies, with $r=\rank_{\R}M$,
\[
 \Dcc(M)\le\lceil\log_2(r+1)\rceil+1,
 \qquad
 \rect(M)\ge\frac1{4(r+1)}.
\]
Moreover $|\cY|\le n+N$.
\end{theorem}

\begin{proof}
For fixed nontrivial $A$ and $a$, the combined $\theta_{n,k}$ gap and
sampling without replacement give failure probability at most
$e^{-\theta_{n,k}N}$.  A union bound over at most $2^n$ sets and $n+2$
relevant consecutive pairs shows that, outside probability $\delta$, every
admissible $A$ is empty, full, a singleton, or a co-singleton.

Because $\cY$ contains the degree-one-unisolvent anchor $\cY_0$, every row
relation lifts to the complete slice and the ambient matrix has the same real
rank $r$.  Every coordinate is nonconstant on $\cY_0$; hence the two-value
condition forces the corresponding rows to be constants, coordinates, or
coordinate complements.  The rank-pooled depth-one protocol from
Theorem~\ref{thm:quadratic-column-hls} gives the communication and rectangle
bounds.
\end{proof}

\begin{theorem}[Deletion-robust linear sparse HLS]
\label{thm:robust-linear-sparse-hls}
If
\[
 N\ge
 \frac8{\theta_{n,k}}
 \left(n\log2+\log(n+2)+\log(1/\delta)\right),
\]
then, with probability at least $1-\delta$, every nontrivial $A$ and every
consecutive pair $\{a,a+1\}$ have at least
\[
 \theta_{n,k}N/2
\]
witnesses $B\in\cY_1$ with $|A\cap B|\notin\{a,a+1\}$.  Consequently all
conclusions of Theorem~\ref{thm:linear-sparse-hls} remain true after an
adversary deletes any $D\subseteq\cY_1$ satisfying
\[
 |D|\le\theta_{n,k}N/4,
\]
while the anchor $\cY_0$ is retained.
\end{theorem}

\begin{proof}
For fixed $A,a$, the number of witnesses in the without-replacement sample
is hypergeometric with mean at least $\theta_{n,k}N$.  Its lower-tail
Chernoff bound gives
\[
 \Pr[X<\theta_{n,k}N/2]\le e^{-\theta_{n,k}N/8}.
\]
Union bound over all $A,a$.  After the stated deletion, at least
$\theta_{n,k}N/4$ witnesses remain, and the rank anchor is unchanged.
\end{proof}

We now remove the requirement that the deterministic rank anchor be
protected.  The first input is an elementary minimum-support estimate for
affine functions on a slice.  Put
\[
 \alpha_{n,k}=\frac{\binom{n-2}{k-1}}{\binom nk}
 =\frac{k(n-k)}{n(n-1)}.
\]

\begin{lemma}[Affine slice support gap]
\label{lem:affine-slice-support-gap}
If a degree-at-most-one function $f$ is not identically zero on
$\binom{[n]}k$, then
\[
 \bigl|\{B\in\tbinom{[n]}k:f(B)\ne0\}\bigr|
 \ge\binom{n-2}{k-1}
 =\alpha_{n,k}\binom nk.
\]
\end{lemma}

\begin{proof}
Write $f(B)=c+\sum_{i\in B}w_i$.  If all $w_i$ are equal, then $f$ is a
nonzero constant on the slice and the claim is immediate.  Otherwise choose
$i,j$ with $w_i\ne w_j$.  For every
$U\in\binom{[n]\setminus\{i,j\}}{k-1}$, pair the two columns
$U\cup\{i\}$ and $U\cup\{j\}$.  Their $f$-values differ by $w_i-w_j$, so at
least one is nonzero.  The pairs are disjoint and there are
$\binom{n-2}{k-1}$ of them.
\end{proof}

We use the following standard finite-population form of VC uniform
convergence.  There is an absolute constant $C_{\rm VC}$ such that, for a
finite set $\Omega$, a class $\mathcal R\subseteq2^\Omega$ of VC dimension at
most $d$, and a uniform $N$-element subset $Y$, the inequalities
\[
 N\ge C_{\rm VC}\varepsilon^{-2}
       \bigl(d\log(2/\varepsilon)+\log(2/\eta)\bigr)
\]
imply, with probability at least $1-\eta$,
\[
 \sup_{R\in\mathcal R}
 \left|\frac{|R\cap Y|}{N}-\frac{|R|}{|\Omega|}\right|
 \le\varepsilon.
\]
This is the classical VC uniform-convergence theorem, with the usual
without-replacement concentration comparison
\cite{VapnikChervonenkis,BardenetMaillard}.  For completeness, if $W$ is a
$d$-dimensional function space, then the support ranges
$\{\operatorname{supp}(f):0\ne f\in W\}$ have VC dimension at most $d$:
shattering $d+1$ points would produce $d+1$ functions whose restrictions are
nonzero multiples of the coordinate basis vectors, contradicting
$\dim W=d$.

\begin{theorem}[Anchor-free fully deletion-robust sparse HLS]
\label{thm:anchor-free-robust-sparse-hls}
Let $\cY$ be a uniformly random $N$-element subset of
$\binom{[n]}k$.  If
\[
 N\ge\max\left\{
 \frac8{\theta_{n,k}}
 \left(n\log2+\log(n+2)+\log(2/\delta)\right),
 \frac{4C_{\rm VC}}{\alpha_{n,k}^2}
 \left(n\log(4/\alpha_{n,k})+\log(4/\delta)\right)
 \right\},
\]
then, with probability at least $1-\delta$, the following holds
simultaneously for every deletion set
\[
 D\subseteq\cY,
 \qquad |D|\le\frac{N}{4}
 \min\{\theta_{n,k},\alpha_{n,k}\}.
\]
For every $\cA\subseteq2^{[n]}$ and integer $a$ satisfying
\[
 |A\cap B|\in\{a,a+1\}
 \qquad(A\in\cA,\ B\in\cY\setminus D),
\]
the matrix $M_{A,B}=|A\cap B|-a$, restricted to columns
$\cY\setminus D$, satisfies, with $r=\rank_{\R}M$,
\[
 \Dcc(M)\le\lceil\log_2(r+1)\rceil+1,
 \qquad
 \rect(M)\ge\frac1{4(r+1)}.
\]
No distinguished subset of columns is protected from deletion.
\end{theorem}

\begin{proof}
Apply the preceding VC estimate with error $\alpha_{n,k}/2$ to the support
ranges of the $n$-dimensional degree-at-most-one slice space.  By
Lemma~\ref{lem:affine-slice-support-gap}, except with probability
$\delta/2$, every nonzero affine slice function is nonzero on at least
$\alpha_{n,k}N/2$ sampled columns.

The hypergeometric lower-tail argument in
Theorem~\ref{thm:robust-linear-sparse-hls}, now with failure probability
$\delta/2$, says that every nontrivial $A$ and consecutive pair have at least
$\theta_{n,k}N/2$ witness columns.  Both events therefore hold with
probability at least $1-\delta$.

After any allowed deletion, every nonzero affine function remains nonzero on
some retained column.  Hence $\cY\setminus D$ is degree-one unisolvent.
Likewise every nontrivial set $A$ retains a witness outside every consecutive
pair, so every admissible row set is empty, full, a singleton, or a
co-singleton.  Unisolvence lifts all row relations to the complete slice and
preserves the actual real rank.  The rank-pooled depth-one protocol used in
Theorem~\ref{thm:quadratic-column-hls} finishes the proof.
\end{proof}

If $\beta n\le k\le(1-\beta)n$, then
\[
 \theta_{n,k}\ge\xi_{n,k}\ge2\beta^2/5.
\]
Thus for fixed $\delta<1$, Theorem~\ref{thm:linear-sparse-hls} proves the
deterministic existence of universal sparse HLS families with only
$O_\beta(n)$ columns.
Theorem~\ref{thm:robust-linear-sparse-hls} uses
$O(\beta^{-2}(n+\log(1/\delta)))$ detector columns and tolerates deletion of
a $\beta^2/10$-fraction of them.
Moreover
\[
 \alpha_{n,k}\ge\beta(1-\beta)\ge2\beta^2/5,
\]
so Theorem~\ref{thm:anchor-free-robust-sparse-hls} also uses
$O(\beta^{-2}(n\log(1/\beta)+\log(1/\delta)))$ columns and tolerates deletion
of a $\beta^2/10$-fraction of the entire column family, with no protected
anchor.

\subsection{Deterministic spectral and design columns}

The preceding existence theorem can be replaced by a deterministic spectral
hypothesis.  For \(B\in Y\subseteq\binom{[n]}k\), write \(x_B\in\{0,1\}^n\)
for its incidence vector and define
\[
 \overline x_Y=\frac1{|Y|}\sum_{B\in Y}x_B,
 \qquad
 \Sigma_Y=\frac1{|Y|}\sum_{B\in Y}
 (x_B-\overline x_Y)(x_B-\overline x_Y)^\top .
\]
Since every column has weight \(k\), one has
\(\Sigma_Y\mathbf1=0\).

\begin{theorem}[Deterministic covariance-gap sparse HLS]
\label{thm:deterministic-spectral-sparse-hls}
Assume \(n\ge4\), \(2\le k\le n-2\), \(Y\ne\varnothing\), and, for every
\(v\perp\mathbf1\),
\begin{equation}\label{eq:deterministic-covariance-gap}
 v^\top\Sigma_Yv\ge\lambda\|v\|_2^2,
 \qquad
 \lambda>\frac{n}{8(n-2)}.
\end{equation}
Then the following holds simultaneously for every
\(\cA\subseteq2^{[n]}\) and integer \(a\) satisfying
\[
 |A\cap B|\in\{a,a+1\}
 \qquad(A\in\cA,\ B\in Y).
\]
For \(M_{A,B}=|A\cap B|-a\) and \(r=\rank_{\R}M\),
\[
 \Dcc(M)\le\lceil\log_2(r+1)\rceil+1,
 \qquad
 \rect(M)\ge\frac1{4(r+1)}.
\]
In particular, this is a deterministic theorem for every prescribed column
family satisfying \eqref{eq:deterministic-covariance-gap}.
\end{theorem}

\begin{proof}
Fix \(A\subseteq[n]\), put \(m=|A|\), and let
\[
 v_A=\mathbf1_A-\frac mn\mathbf1\in\mathbf1^\perp.
\]
Because \(\Sigma_Y\mathbf1=0\),
\[
 \Var_{B\in Y}|A\cap B|
 =\mathbf1_A^\top\Sigma_Y\mathbf1_A
 =v_A^\top\Sigma_Yv_A
 \ge\lambda\frac{m(n-m)}n.
\]
If \(2\le m\le n-2\), then
\[
 \frac{m(n-m)}n\ge\frac{2(n-2)}n,
\]
so \eqref{eq:deterministic-covariance-gap} makes the variance strictly larger
than \(1/4\).  This contradicts the fact that a random variable supported on
\(\{a,a+1\}\) has variance at most \(1/4\).  Hence every admissible \(A\) is
empty, full, a singleton, or a co-singleton.

The same spectral gap proves degree-one unisolvence without a separate
anchor.  Indeed, suppose
\[
 f(B)=c+w^\top x_B
\]
vanishes on \(Y\).  Writing
\(w=w_0+t\mathbf1\) with \(w_0\perp\mathbf1\), the empirical variance of
\(f\) is zero, whereas \eqref{eq:deterministic-covariance-gap} gives
\[
 0=w_0^\top\Sigma_Yw_0\ge\lambda\|w_0\|_2^2.
\]
Thus \(w_0=0\), and \(f\) is constant on the complete weight-\(k\) slice.
Its value on the nonempty set \(Y\) is zero, so it vanishes on the complete
slice.

Consequently every row relation on \(Y\) lifts to the complete slice and the
ambient and restricted matrices have the same real rank.  The positive gap
also makes every coordinate nonconstant on \(Y\).  Hence the two-value
condition forces singleton rows to be coordinates, co-singleton rows to be
coordinate complements, and the remaining classified rows to be constants,
all with Boolean extensions to the complete slice.  The rank-pooled
depth-one protocol from Theorem~\ref{thm:quadratic-column-hls} gives the
communication bound, and a largest transcript rectangle gives the stated
density.
\end{proof}

A large, standard source of deterministic spectral families is furnished by
block designs.

\begin{corollary}[\(2\)-design columns]
\label{cor:two-design-spectral-hls}
Let \(Y\) be the block family of a simple \(2\)-\((n,k,\lambda_2)\) design.
If
\begin{equation}\label{eq:design-spectral-threshold}
 8k(n-k)(n-2)>n^2(n-1),
\end{equation}
then all conclusions of
Theorem~\ref{thm:deterministic-spectral-sparse-hls} hold.
\end{corollary}

\begin{proof}
Write \(b=|Y|\) and let \(r_2\) be the number of blocks through one point.
The design identities give
\[
 \frac{r_2}{b}=\frac kn,
 \qquad
 \frac{\lambda_2}{b}=\frac{k(k-1)}{n(n-1)}.
\]
It follows entrywise that
\[
 \Sigma_Y=\frac{k(n-k)}{n(n-1)}
 \left(I-\frac1nJ\right).
\]
Thus the eigenvalue on \(\mathbf1^\perp\) is
\(k(n-k)/[n(n-1)]\), and
\eqref{eq:design-spectral-threshold} is exactly the strict threshold in
\eqref{eq:deterministic-covariance-gap}.
\end{proof}

The design corollary yields an explicit information-theoretically minimal
family on an infinite sequence of central layers.

\begin{corollary}[Explicit Sylvester--Hadamard columns]
\label{cor:sylvester-hadamard-sparse-hls}
For every \(t\ge3\), put
\[
 n=2^t-1,
 \qquad
 k=2^{t-1}-1.
\]
There is an explicit family
\[
 Y_t\subseteq\binom{[n]}k,
 \qquad |Y_t|=n,
\]
for which all conclusions of
Theorem~\ref{thm:deterministic-spectral-sparse-hls} hold simultaneously for
every adjacent-two-value row family.  The column count \(n\) is minimal for
degree-one unisolvence.
\end{corollary}

\begin{proof}
Let \(H_t\) be the Sylvester Hadamard matrix, recursively defined by
\[
 H_1=\begin{pmatrix}1&1\\1&-1\end{pmatrix},
 \qquad
 H_{t+1}=\begin{pmatrix}H_t&H_t\\H_t&-H_t\end{pmatrix}.
\]
It has order \(4u=2^t\), where \(u=2^{t-2}\), and its first row and column are
all \(+1\).  Delete that row and column.  For every remaining row, take the
set of positions carrying \(+1\).  Orthogonality shows that every block has
size \(2u-1\), and every two blocks meet in exactly \(u-1\) points.  Hence
these \(4u-1=n\) blocks form a symmetric
\(2\)-\((4u-1,2u-1,u-1)\) design.

Its covariance eigenvalue on \(\mathbf1^\perp\) is
\[
 \frac{k(n-k)}{n(n-1)}=\frac{u}{4u-1}.
\]
For \(u\ge2\),
\[
 \frac{u}{4u-1}>\frac{4u-1}{8(4u-3)}
\]
because the difference after clearing positive denominators is
\(16u^2-16u-1>0\).  Theorem~\ref{thm:deterministic-spectral-sparse-hls}
therefore applies.  Finally, the degree-at-most-one slice space has dimension
\(n\), so an unisolvent family cannot have fewer than \(n\) columns.
\end{proof}

There is an analogous minimal sparse family for every bounded-alphabet
complete product.  For
$\Omega_{\rm prod}=\prod_i[m_i]$, decompose each one-coordinate function
space into constants and its $(m_i-1)$-dimensional mean-zero part.  The
degree-at-most-$s$ product space then has dimension
\begin{equation}\label{eq:product-degree-dimension}
 d_s^{\rm prod}
 =\sum_{\substack{J\subseteq[n]\\|J|\le s}}
   \prod_{i\in J}(m_i-1).
\end{equation}
It is invariant under the transitive action of $\prod_iS_{m_i}$.

\begin{corollary}[Minimal sparse bounded-product columns]
\label{cor:dpp-product}
Sample $Y$ from the rank-$d_s^{\rm prod}$ projection DPP of the
degree-at-most-$s$ space on $\Omega_{\rm prod}$.  With probability one,
\[
 |Y|=d_s^{\rm prod}
\]
and every Boolean degree-$s$ row family restricted to $Y$ satisfies the
communication bound in Theorem~\ref{thm:low-degree-product-domain}, with its
actual restricted real rank.
\end{corollary}

\begin{proof}
Apply Theorems~\ref{thm:dpp-unisolvent} and
\ref{thm:unisolvent-transfer}, using
\eqref{eq:product-degree-dimension}.
\end{proof}

For fixed $Q,s$,
\[
 d_s^{\rm prod}
 \le\sum_{j=0}^s\binom nj(Q-1)^j
 =O_{Q,s}(n^s),
\]
whereas $|\Omega_{\rm prod}|\ge2^n$.  Thus this construction is again
information-theoretically minimal and exponentially sparse.

For fixed $s$, the uniform Johnson construction uses $O_s(n^s\log n)$
columns.  A middle Johnson layer has $\Theta(2^n/\sqrt n)$ columns, so the retained family is
exponentially sparse.  Taking any fixed $\delta<1$ also proves deterministic
existence of such sparse families.  More basically, every $d$-dimensional
function space has some unisolvent set of exactly $d$ points, by choosing a
basis of its evaluation functionals.  Projection-DPP sampling attains this
minimum exactly, while the uniform theorem gives an elementary
without-replacement construction at logarithmic oversampling.

\section{The product of Johnson middle layers}

Let
\[
 [d]=C_1\sqcup\cdots\sqcup C_m,
 \qquad |C_i|=n,
\]
where $n\ge4$ is even, and define
\begin{equation}\label{eq:product-family}
 \cY_{n,m}=\{B\subseteq[d]:|B\cap C_i|=n/2\text{ for all }i\}.
\end{equation}

\begin{theorem}[Product-Johnson restricted Log-Rank theorem]
\label{thm:product}
Let $\cA\subseteq2^{[d]}$ be nonempty and suppose
\[
 |A\cap B|\in\{a,a+1\}
 \qquad(A\in\cA,\ B\in\cY_{n,m}).
\]
Then the matrix \eqref{eq:matrix} satisfies
\[
 \rank_{\R}M\le d+1,
 \qquad
 \rect(M)\ge\frac1{4(d+1)},
 \qquad
 \Dcc(M)\le\lceil\log_2(\rank_{\R}M+1)\rceil+4.
\]
\end{theorem}

\begin{proof}
Choose $B$ uniformly from $\cY_{n,m}$.  The restrictions $B\cap C_i$ are
independent uniform middle-layer sets.  Put $s_i=|A\cap C_i|$.  The
hypergeometric variance formula gives
\begin{equation}\label{eq:block-var}
 \Var|A\cap B|
 =\sum_{i=1}^m\frac{s_i(n-s_i)}{4(n-1)}.
\end{equation}
By Lemma~\ref{lem:two-point-var}, the left side is at most $1/4$.

Let
\[
 h(A)=\sum_i\min\{s_i,n-s_i\},
\]
the Hamming distance from $A$ to the nearest union of blocks.  If
$h(A)\ge2$, then either two errors occur in one block or errors occur in at
least two blocks.  In the first case the corresponding numerator in
\eqref{eq:block-var} is at least $2(n-2)$; in the second it is at least
$2(n-1)$.  Here we use that $r(n-r)$ is increasing for
$0\le r\le n/2$, so additional error units can only increase the relevant
lower bound.  Since $n\ge4$, both cases make \eqref{eq:block-var} strictly
larger than $1/4$, a contradiction.  Therefore $h(A)\le1$.

Write
\[
 A=C_J\triangle E,
 \qquad C_J=\bigcup_{i\in J}C_i,
 \qquad E\in\{\varnothing,\{1\},\ldots,\{d\}\}.
\]
Pigeonhole the rows by $E$, retaining at least a $1/(d+1)$ fraction.  If
$E=\varnothing$, every retained row has constant intersection with every
column because each block count is fixed.  If $E=\{e\}$, additionally fix
whether $e\in B$; exactly half the columns lie in each class.  Every retained
row is now constant on the retained columns.  Keep the larger of the rows
with constant value $a$ and those with constant value $a+1$.  The total loss
is at most $4(d+1)$.

The rank assertion follows from \eqref{eq:rank-factorization}.
\end{proof}

\subsection{A rank-sensitive one-way protocol}

The row classification in the proof gives a full protocol, not only one
large rectangle.  If $E=\varnothing$, the row is constant.  If
$E=\{e\}$ and the block containing $e$ is not in $J$, then
\[
 |A\cap B|=|J|n/2+\1_{\{e\in B\}},
\]
so the nonconstant matrix row is exactly the coordinate function
$B\mapsto\1_{\{e\in B\}}$.  If that block is in $J$, then
\[
 |A\cap B|=|J|n/2-\1_{\{e\in B\}},
\]
and the nonconstant row is the complement of the coordinate function.

Let $S\subseteq[d]$ be the coordinates occurring in nonconstant rows and
put $k=|S|$.  Alice can send one of two constant messages, or the index of
$e\in S$ together with the complement bit.  Thus there is a one-way protocol
with at most $2k+2$ messages.

We next relate $k$ to the actual real rank $r=\rank_{\R}M$.  Let $X_S$ be
the matrix whose rows are the coordinate functions
$B\mapsto\1_{\{e\in B\}}$, $e\in S$.  If $f$ blocks are fully contained in
$S$, then
\begin{equation}\label{eq:coordinate-rank}
 \rank X_S=
 \begin{cases}
 k,&f=0,\\
 k-f+1,&f>0.
 \end{cases}
\end{equation}
To see this, vary one middle-layer block while fixing all other blocks.  A
linear combination of coordinate functions is constant on that block only
when all its coefficients are equal, by exchanging one selected and one
unselected coordinate.  A partially represented block therefore contributes
no relation.  Each fully represented block contributes its fixed-sum
relation, and the $f$ block sums span one common constant direction.

Since $f\le k/n$,
\begin{equation}\label{eq:k-rankx}
 \rank X_S\ge(1-1/n)k.
\end{equation}
Choose one nonconstant row of $M$ for every $e\in S$.  The resulting matrix
is obtained from $X_S$ by invertible row sign changes and a rank-one constant
perturbation.  Hence
\[
 r\ge\rank X_S-1\ge(1-1/n)k-1,
\]
and, because $n\ge4$,
\begin{equation}\label{eq:k-rank}
 k\le\frac n{n-1}(r+1)\le\frac43(r+1).
\end{equation}
Encoding the $2k+2$ possible one-way messages costs at most
$\lceil\log_2(r+1)\rceil+3$ bits, which implies the stated bound.

\begin{corollary}[Optimality of the logarithmic order]
\label{cor:optimal}
For fixed even $n\ge4$ and growing $m$, the order
$\Dcc(M)=\Theta(\log\rank M)$ in Theorem~\ref{thm:product} is best possible.
\end{corollary}

\begin{proof}
Take $\cA=\{\{e\}:e\in[d]\}$ and $a=0$.  Then $M$ is the coordinate
incidence matrix of $\cY_{n,m}$.  Formula \eqref{eq:coordinate-rank} with
all $m$ blocks fully represented gives
\[
 \rank_{\R}M=d-m+1.
\]
Every deterministic $c$-bit protocol partitions a Boolean matrix into at
most $2^c$ monochromatic rectangles, each of real rank at most one.  Hence
$c\ge\log_2\rank M$.  Theorem~\ref{thm:product} supplies the matching upper
bound up to an additive constant.
\end{proof}

The same proof gives a parameterized theorem for arbitrary products of
Johnson layers.  Let the blocks have sizes $n_i\ge2$, choose
$1\le k_i\le n_i-1$, and put
\[
 \cY(\mathbf n,\mathbf k)
 =\{B\subseteq[d]:|B\cap C_i|=k_i\text{ for all }i\}.
\]
For $0\le s_i\le n_i$, define
\[
 V(\mathbf s)=
 \sum_i
 \frac{s_i(n_i-s_i)k_i(n_i-k_i)}{n_i^2(n_i-1)},
 \qquad
 h(\mathbf s)=\sum_i\min\{s_i,n_i-s_i\},
\]
and the exact weighted tube radius
\begin{equation}\label{eq:product-tube}
 t(\mathbf n,\mathbf k)
 =\max\{h(\mathbf s):V(\mathbf s)\le1/4\}.
\end{equation}

\begin{theorem}[General product-Johnson theorem]
\label{thm:general-product}
If \eqref{eq:two-level} holds with
$\cB=\cY(\mathbf n,\mathbf k)$ and
$t=t(\mathbf n,\mathbf k)$, then
\[
 \rank_{\R}M\le d+1,
 \qquad
 \rect(M)\ge
 \frac1{2^{t+1}\sum_{j=0}^t\binom dj}.
\]
\end{theorem}

\begin{proof}
The block restrictions are independent and the hypergeometric variance is
exactly $V(\mathbf s)$.  Lemma~\ref{lem:two-point-var} therefore forces
$h(\mathbf s)\le t$.  Write each row as a block union symmetric-differenced
with an error set of size at most $t$.  Pigeonhole the error set, then fix its
exact trace on the right; this costs at most
$2^t\sum_{j\le t}\binom dj$.  Every retained row is now constant across the
retained columns, and a final two-colour pigeonhole completes the proof.
\end{proof}

The parameterized theorem also gives a communication protocol.

\begin{theorem}[Protocol for general product layers]
\label{thm:general-product-protocol}
Let $r=\rank_{\R}M$ and $t=t(\mathbf n,\mathbf k)$.  If $t=0$, then $M$ has
only constant rows.  For every $t\ge1$,
\[
 \Dcc(M)=O\bigl(t\log(r+2)+\log(t+2)\bigr).
\]
In particular, if $t\le(\log(r+2))^C$, then the Log-Rank Conjecture holds for
this product-layer matrix.
\end{theorem}

\begin{proof}
Every row has a sparse affine representation on the column family,
\begin{equation}\label{eq:sparse-row-rep}
 M_A(B)=q_A+\sum_{e\in E_A}\sigma_{A,e}\1_{\{e\in B\}},
 \qquad |E_A|\le t,
\end{equation}
where $\sigma_{A,e}\in\{-1,1\}$ and $q_A\in\Z$.  The two-level condition
forces $q_A\in[-t,t+1]$.

Choose $r$ matrix rows spanning the row space, and let $S_0$ be the union of
their error supports.  Then $|S_0|\le rt$.  We claim that the union $S$ of
the error supports of all rows satisfies
\begin{equation}\label{eq:used-coordinate-bound}
 |S|\le2|S_0|\le2rt.
\end{equation}

To prove the claim, consider the evaluation map from affine coordinate forms
to functions on $\cY(\mathbf n,\mathbf k)$.  Its kernel consists exactly of
forms whose coefficients are constant on each block, with the constant term
chosen to cancel $\sum_i\alpha_i k_i$.  Indeed, varying a single block and
exchanging one selected and one unselected coordinate forces all
coefficients in that block to be equal.

For an arbitrary row, subtract a linear combination of the chosen basis-row
representations.  The difference lies in this kernel.  If a block contains a
new error coordinate outside $S_0$, its kernel coefficient is nonzero;
therefore every coordinate of that block outside $S_0$ must belong to the
row's error set.  Use the canonical nearest-block-union representation, so
every block contributes at most half of its coordinates to an error set.
Consequently such a block has at most as many coordinates outside $S_0$ as
inside $S_0$.  Summing over the disjoint blocks that contain new coordinates
proves \eqref{eq:used-coordinate-bound}.

Alice now sends $q_A$, the subset $E_A\subseteq S$, and its signs.  The
number of possible messages is at most
\[
 (2t+2)\sum_{j=0}^t\binom{|S|}{j}2^j.
\]
Using $|S|\le2rt$ and the standard binomial-sum estimate gives the claimed
communication bound.  Bob evaluates \eqref{eq:sparse-row-rep} from his
column.
\end{proof}

\subsection{All interior product layers}

The strongest product theorem does not require a variance-radius estimate.

\begin{theorem}[Interior product-Johnson Log-Rank theorem]
\label{thm:interior-product}
Assume
\[
 2\le k_i\le n_i-2
 \qquad\text{for every block }i.
\]
If \eqref{eq:two-level} holds with
$\cB=\cY(\mathbf n,\mathbf k)$, then
\[
 \rect(M)\ge\frac1{4(d+1)},
 \qquad
 \Dcc(M)\le\lceil\log_2(\rank_{\R}M+1)\rceil+4.
\]
The logarithmic order is optimal up to an additive constant.
\end{theorem}

\begin{proof}
Fix a row $A$ and put $s_i=|A\cap C_i|$.  As $B_i$ ranges over all
$k_i$-subsets of $C_i$, the possible values of $|A\cap B_i|$ are all integers
in the interval
\begin{equation}\label{eq:local-range}
 \left[
  \max\{0,k_i-(n_i-s_i)\},\ \min\{k_i,s_i\}
 \right].
\end{equation}
The choices in distinct blocks are independent, so the range of
$|A\cap B|$ is the sum of these integer intervals.  Since the total range has
size at most two, the sum of their widths is at most one.

For $2\le k_i\le n_i-2$, the interval in \eqref{eq:local-range} has width zero exactly when
$s_i\in\{0,n_i\}$, and has width one only when
$s_i\in\{1,n_i-1\}$.  Indeed, if $2\le s_i\le n_i-2$ and
$s_i+k_i\le n_i$, its lower endpoint is zero and its upper endpoint is at
least two; if $s_i+k_i>n_i$, its width is
$n_i-\max\{s_i,k_i\}\ge2$.

Consequently all blocks are empty or full in $A$, except possibly one block
where $A$ differs from empty or full in one coordinate.  Thus every row is
constant, a coordinate function, or its complement.  The rectangle proof of
Theorem~\ref{thm:product} applies verbatim.

For the rank-sensitive protocol, let $S$ be the set of queried coordinates.
The coordinate-rank calculation \eqref{eq:coordinate-rank} remains valid for
unequal blocks.  Since every block has size at least four, the number of fully
represented blocks is at most $|S|/4$, and hence
$\rank X_S\ge3|S|/4$.  The rank-one perturbation argument then gives
$|S|\le4(r+1)/3$, completing the protocol proof.  The singleton-row example
again proves optimality.
\end{proof}

In the language of Theorem~\ref{thm:atomic-transfer}, interior Johnson factors
have atom-growth exponent $\alpha=1$ and constant $C\le4/3$: the centered
rank of $k$ queried coordinates is $k$ unless the entire block is queried,
in which case it is $k-1\ge3k/4$.

\begin{proposition}[Sharp boundary obstruction]\label{prop:boundary}
If blocks of weight one are allowed, the class contains every Boolean matrix
with the corresponding number of columns.
\end{proposition}

\begin{proof}
Take one block $C$ of size $N$ with $k=1$, so columns are the singletons
$\{j\}$, $j\in C$.  Given any Boolean matrix with these $N$ columns, identify
each row with the subset $A\subseteq C$ of its one-entries.  Then
\[
 |A\cap\{j\}|\in\{0,1\}
\]
and the resulting adjacent-two-value matrix is the original matrix.  The case
$k=N-1$ follows by complementation.
\end{proof}

\section{Products of interior multislices}

The product/rank argument is not specific to two-color slices.  For a vector
of positive integers $\boldsymbol\kappa=(\kappa_1,\ldots,\kappa_\ell)$ with
sum $N$, define the multislice
\[
 \Omega(\boldsymbol\kappa)
 =\left\{x\in[\ell]^N:
   |\{p:x_p=c\}|=\kappa_c\text{ for every }c\right\}.
\]
A function has degree at most one if it lies in the affine span of the
coordinate-color indicators $1_{\{x_p=c\}}$.

Consider a product of multislices
$\Omega=\prod_{i=1}^m\Omega(\boldsymbol\kappa^{(i)})$, each with alphabet
size at most $Q$.  A Boolean matrix with column set $\Omega$ is called
row-degree-one if every row is a Boolean degree-one function on this product.

\begin{theorem}[Interior multislice product theorem]
\label{thm:multislice}
Assume every color multiplicity in every factor is at least two.  If $M$ is a
row-degree-one Boolean matrix on $\Omega$, then
\[
 \Dcc(M)\le
 \left\lceil\log_2(\rank_{\R}M+1)\right\rceil+Q+4.
\]
For fixed alphabet bound $Q$, this is optimal-order Log-Rank.
\end{theorem}

\begin{proof}
Filmus--Ihringer show that a Boolean degree-one function on an interior
multislice depends only on the color of one position
\cite[Lemma~8.1]{FilmusIhringer}.  A Boolean degree-one function on a product
is a sum of one-factor functions.  Two nonconstant factors would produce at
least three sums, so every nonconstant row depends on one position in one
factor and has the form
\[
 x\longmapsto1_{\{x_p\in J\}}
\]
or its complement, for a nontrivial color subset $J$.  Up to complement there
are at most $2^{Q-1}-1$ atoms per position.

Let $P$ be the number of positions used by the matrix rows.  We claim
\begin{equation}\label{eq:multislice-position-rank}
 P\le2(\rank_{\R}M+1).
\end{equation}
Choose one atom for each used position.  Within one multislice factor, an
additive coordinate form that vanishes identically has the form
$g_p(c)=a_p+h_c$, with its constants satisfying the single histogram
relation.  This follows by exchanging the colors at two positions.  Hence a
proper subset of positions contributes independent nonconstant directions;
if all $N$ positions are used, the centered rank is at least
$N-\ell\ge N/2$, since every color multiplicity is at least two.  Across
different factors the centered function spaces form a direct sum.  Therefore
the centered rank of one selected atom per used position is at least $P/2$.

Replacing complements by positive atoms enlarges the row span by at most the
constant function, so \eqref{eq:multislice-position-rank} follows.  The total
number of used atoms is at most $2^{Q-1}P\le2^Q(r+1)$.  Alice sends the atom
index and complement bit, proving the stated protocol bound.
\end{proof}

Thus interior multislice factors satisfy the atomic growth hypothesis with
$\alpha=1$ and $C\le2^Q$.

\begin{proposition}[Sharp multislice boundary]
If a color of multiplicity one is allowed, row-degree-one matrices on a
single multislice include every Boolean matrix with $N$ columns.
\end{proposition}

\begin{proof}
Use the binary multislice with histogram $(1,N-1)$.  A column is identified
by the unique position carrying the rare color.  For any subset
$J\subseteq[N]$, the function
\[
 x\longmapsto\sum_{p\in J}1_{\{x_p=\text{rare}\}}
\]
is Boolean and has degree one; these functions realize arbitrary row
supports.
\end{proof}

\section{Sparsity and separation from affine sections}

The family \eqref{eq:product-family} lies in the global middle layer and has
relative density
\begin{equation}\label{eq:relative-density}
 \frac{|\cY_{n,m}|}{\binom d{d/2}}
 =\frac{\binom n{n/2}^m}{\binom{mn}{mn/2}}.
\end{equation}
For fixed $n$ and $m\to\infty$, Stirling's formula shows that
\eqref{eq:relative-density} is $\exp(-\Theta_n(d))$.  For $n=4$ it is
\[
 \Theta\bigl(\sqrt m(3/8)^m\bigr).
\]

The family is also generally not a complete affine-code section.  For two
blocks of size four, label each block by $\{1,2,3,4\}$.  In the first block
take
\[
 \{1,2\},\quad\{1,3\},\quad\{2,3\},
\]
whose XOR is empty.  In the second take
\[
 \{1,2\},\quad\{1,3\},\quad\{1,4\},
\]
whose XOR is the full block.  Pairing corresponding choices gives three
product columns with triple XOR of block weights $(0,4)$.  This XOR is still
a global middle-layer set but is not in the product family.  Any binary
affine coset containing the three columns contains their triple XOR, so its
intersection with the middle layer cannot equal $\cY_{4,2}$.

\section{Boolean tubes around a defect space}

Let $D\le\R^d$.  Define its Boolean templates
\[
 \cC(D)=\{C\subseteq[d]:\1_C\in D\}
\]
and the Hamming distance to that family
\[
 h_D(A)=\min_{C\in\cC(D)}|A\triangle C|.
\]
For $\lambda>0$, define the exact Boolean tube radius
\begin{equation}\label{eq:tube-radius}
 \tau_\lambda(D)=
 \max\left\{h_D(A):
  \dist(\1_A,D)^2\le\frac1{4\lambda}\right\}.
\end{equation}
The maximum exists because the Boolean cube is finite and $\varnothing$ is a
template.

A convenient sufficient estimate is a Boolean Hamming margin
\begin{equation}\label{eq:boolean-margin}
 \dist(\1_A,D)^2\ge\gamma h_D(A),
\end{equation}
which implies
\[
 \tau_\lambda(D)\le\left\lfloor(4\lambda\gamma)^{-1}\right\rfloor.
\]
The exact radius is often sharper than the global margin.  For the span of
one all-ones vector it recovers the distinction between singleton and
two-coordinate errors that is lost in a uniform factor-$1/2$ estimate.

\section{Integer generator entropy}

Suppose nonzero $z_1,\ldots,z_s\in\Z^d$ span $D$ over $\R$.  Put
\[
 e_i=\|z_i\|_2^2,
 \qquad h_i=\|z_i\|_1.
\]
For $u\ge0$ and an integer $h\ge1$, define
\[
 \phi(u,h)=
 \begin{cases}
 H_2(4u)+4u\log h,&u\le1/8,\\
 \log(h+1),&u>1/8,
 \end{cases}
\]
where $H_2$ is binary entropy in natural units.

\begin{lemma}[Integer entropy envelope]\label{lem:entropy-envelope}
Let $Z$ be integer-valued, supported on at most $h+1$ consecutive integer
values, and suppose $\Var Z\le u$.  Then $H(Z)\le\phi(u,h)$.
\end{lemma}

\begin{proof}
Let $p$ be the largest atom.  For an independent copy $Z'$, the identity
\[
 \Var Z=\frac12\mathbb E(Z-Z')^2
\]
shows that $\Var Z\ge p(1-p)$.  Moreover, if $p\le1/2$, then
\[
 \Var Z\ge\frac12\Pr[Z\ne Z']
 =\frac12\left(1-\sum_x\Pr[Z=x]^2\right)\ge1/4.
\]
Consequently, when $u\le1/8$, one has $p>1/2$ and
$1-p\le2u\le4u$.  Entropy is maximized by spreading the remaining mass
uniformly on the other $h$ support points, giving
$H(Z)\le H_2(4u)+4u\log h$.  For $u>1/8$, use
$H(Z)\le\log(h+1)$.
\end{proof}

Let $X=\1_B$ for a uniform random column $B\in\cB$, and suppose
\begin{equation}\label{eq:inside-upper}
 v^\mathsf T\Sigma v\le\kappa\|v\|_2^2
 \qquad(v\in D),
\end{equation}
where $\Sigma=\Cov(X)$.  Define
\begin{equation}\label{eq:psi}
 \Psi_Z=\sum_{i=1}^s\phi(\kappa e_i,h_i).
\end{equation}
The integer statistics $T_i=z_i\cdot X$ have variance at most $\kappa e_i$
and range size at most $h_i+1$.  Entropy subadditivity gives
\begin{equation}\label{eq:joint-atom}
 \max_t\Pr[(T_1,\ldots,T_s)=t]\ge e^{-\Psi_Z}.
\end{equation}

\section{The spectral--lattice rectangle theorem}

\begin{theorem}[Spectral--lattice rectangle theorem]
\label{thm:master}
Let $\cA,\cB\subseteq2^{[d]}$ be nonempty and satisfy
\eqref{eq:two-level}.  Let $D\le\R^d$ be spanned by nonzero integer vectors
$z_1,\ldots,z_s$.  Suppose
\begin{equation}\label{eq:outside-gap}
 \Sigma\succeq\lambda P_{D^\perp}
\end{equation}
and \eqref{eq:inside-upper}.  Put $t=\tau_\lambda(D)$ and
$S(d,t)=\sum_{j=0}^t\binom dj$.  Then
\begin{equation}\label{eq:master-bound}
 \rect(M)\ge
 \frac{e^{-\Psi_Z}}{2^{t+1}S(d,t)}.
\end{equation}
\end{theorem}

\begin{proof}
For fixed $A\in\cA$, Lemma~\ref{lem:two-point-var} and
\eqref{eq:outside-gap} give
\[
 \lambda\dist(\1_A,D)^2
 \le\1_A^\mathsf T\Sigma\1_A
 =\Var|A\cap B|
 \le1/4.
\]
By the definition of $t$, choose $C(A)\in\cC(D)$ and $E(A)$ with
\[
 A=C(A)\triangle E(A),
 \qquad |E(A)|\le t.
\]
There are at most $S(d,t)$ error sets.  Retain a largest error fiber.

For the fixed error $E$, refine $\cB$ to a largest atom of
\[
 (z_1\cdot\1_B,\ldots,z_s\cdot\1_B, B\cap E).
\]
By \eqref{eq:joint-atom}, this retains at least
$e^{-\Psi_Z}2^{-t}$ of the columns.  Since every template is a real linear
combination of the generators, the generator statistics fix
$|C(A)\cap B|$.  The exact trace on $E$ and the identity
\[
 |(C\triangle E)\cap B|+2|C\cap E\cap B|
 =|C\cap B|+|E\cap B|
\]
then fix $|A\cap B|$ for every retained row.  Retain the larger row colour,
losing at most two.  Multiplying the losses gives
\eqref{eq:master-bound}.
\end{proof}

\begin{remark}[Dimension-free character]
Theorem~\ref{thm:master} has no explicit factor exponential in $\dim D$ or in
$|\cC(D)|$.  Arithmetic complexity is paid through the entropy of the
generator statistics.
\end{remark}

If $E_2=\sum_i\|z_i\|_2^2$,
$H=\max(1,\max_i\|z_i\|_1)$, and $0<\kappa\le1$, then
\begin{equation}\label{eq:coarse-psi}
 \Psi_Z\le
 O\left(\kappa E_2\log\frac{e(H+1)}{\kappa}\right),
\end{equation}
with the convention that the right side is zero for $\kappa=0$.

Indeed, put $u_i=\kappa e_i$ and $U=\sum_i u_i=\kappa E_2$.
For $u_i\le1/8$, the bound
$H_2(x)\le x\log(e/x)$ gives
\[
 \phi(u_i,h_i)\le4u_i\log\frac{eH}{u_i}.
\]
Since the number of nonzero integer generators is at most $E_2$, concavity
of $x\mapsto x\log(eH/x)$ bounds the sum of these terms by
$O(U\log(eH/\kappa))$.  There are at most $8U$ indices with
$u_i>1/8$, and each contributes at most $\log(H+1)$.  This proves
\eqref{eq:coarse-psi}.

\section{Robust transfer from a nearby rational subspace}

The exact arithmetic hypothesis admits a perturbative extension.  Let
$D_0$ be the true spectral defect space and $D_\mathbb Q$ an
integer-generated subspace, with projections $P_0,P_\mathbb Q$ satisfying
\begin{equation}\label{eq:projector-close}
 \|P_0-P_\mathbb Q\|_{\mathrm{op}}\le\varepsilon.
\end{equation}
Assume $D_0$ has outside gap $\lambda$ and Boolean margin $\gamma_0$, while
$D_\mathbb Q$ has Boolean margin $\gamma_\mathbb Q$.  Suppose covariance on
$D_0$ is at most $\kappa$.

Every Boolean template in $D_0$ lies within Euclidean distance
$\varepsilon\sqrt d$ of $D_\mathbb Q$, hence within Hamming distance
\begin{equation}\label{eq:rational-error}
 s=\left\lfloor\frac{\varepsilon^2d}{\gamma_\mathbb Q}\right\rfloor
\end{equation}
of an exact Boolean template there.  Moreover, since
$\|\Sigma\|_{\mathrm{op}}\le\operatorname{tr}\Sigma\le d/4$, covariance on
$D_\mathbb Q$ is at most
\begin{equation}\label{eq:kappa-transfer}
 \kappa_\mathbb Q=2\kappa+\frac d2\varepsilon^2.
\end{equation}

\begin{theorem}[Robust rational transfer]\label{thm:robust}
Under the preceding assumptions, let
\[
 t=\left\lfloor(4\lambda\gamma_0)^{-1}\right\rfloor,
 \qquad T=t+s.
\]
If $z_i$ are integer generators for $D_\mathbb Q$ and $\Psi_\mathbb Q$ is
defined using $\kappa_\mathbb Q$, then
\[
 \rect(M)\ge
 \frac{e^{-\Psi_\mathbb Q}}
 {2^{T+1}\sum_{j=0}^T\binom dj}.
\]
\end{theorem}

\begin{proof}
The outside gap and margin put every row within $t$ coordinates of a Boolean
template in $D_0$.  Projector closeness and the rational margin put that
template within $s$ coordinates of one in $D_\mathbb Q$.  The Hamming
triangle inequality gives total error at most $T$.  The covariance estimate
\eqref{eq:kappa-transfer} permits the extraction argument of
Theorem~\ref{thm:master} on the rational space with the enlarged error
budget; the outside-gap step has already been supplied by $D_0$.
\end{proof}

For an orthonormal basis $w_1,\ldots,w_q$ of $D_0$ with rational
approximants $r_i=z_i/M_i$ satisfying $\|r_i-w_i\|_2\le\eta$ and
$M_i\le H_0$, one has, when $\eta\sqrt q<1$,
\[
 \|P_0-P_\mathbb Q\|_{\mathrm{op}}
 \le\frac{2\eta\sqrt q}{1-\eta\sqrt q},
 \qquad
 \sum_i\|z_i\|_2^2\le qH_0^2(1+\eta)^2.
\]
Thus the robust theorem gives an explicit endpoint for simultaneous inverse
Littlewood--Offord basis approximation.

\section{Relation to previous work}

The cross-intersection formulation is due to Hambardzumyan, Lovett, and
Shirley \cite{HambardzumyanLovettShirley}.  Their theorem is an equivalence,
not a solution for the product family considered here.  The general
$2^{-O(\sqrt r)}$ rectangle bound of Sudakov--Tomon applies, but is
exponentially weaker than Theorem~\ref{thm:product} when $r\asymp d$
\cite{SudakovTomon}.

Kupavskii--Weltge prove
$|\cA||\cB|\le(d+1)2^d$ for full-span binary scalar products, and
Kupavskii--Tsarev prove sharp near-extremal stability
\cite{KupavskiiWeltge,KupavskiiTsarev}.  Those are absolute cardinality
statements.  They do not yield a relative monochromatic rectangle inside a
given exponentially sparse product family.

Rao--Yehudayoff prove powerful anti-concentration results for large cube
families and certain finite-field structured distributions
\cite{RaoYehudayoff}.  Nguyen proves an inverse Littlewood--Offord theorem for
the complete half layer \cite{NguyenCombinatorial}.  These results concern
ambient concentration and do not directly give Theorem~\ref{thm:master} for
an arbitrary sparse right family.

There is also a substantial Erd\H{o}s--Ko--Rado literature for direct
products of uniform families.  For example, Yao--Lv--Wang determine extremal
$t$-intersecting subfamilies in a multipart setting \cite{YaoLvWang}.  Those
results impose the one-family condition $|A\cap A'|\ge t$.  They do not treat
a cross pair with a prescribed full product on one side, two adjacent
intersection values, or the resulting rank-sensitive communication protocol.

The one-block local classification in
Theorem~\ref{thm:interior-product} is consistent with the Boolean degree-one
theorem of Filmus--Ihringer \cite{FilmusIhringer}.  Their paper does not state
the product communication theorem, the rank-sensitive support argument, or
the weight-one universality boundary.  Accordingly, novelty here rests on
those global and communication-complexity components, not on re-claiming the
single-scheme degree-one classification.

Likewise, Filmus--Ihringer's multislice classification
\cite[Lemma~8.1]{FilmusIhringer} is the local input to
Theorem~\ref{thm:multislice}.  The product assembly, centered-rank estimate,
rank-sensitive protocol, and count-one universality boundary are the claims
made here.

The low-degree slice theorems use published results of Filmus--Ihringer,
Chiarelli--Hatami--Saks, Midrijanis, Filmus, and
Dafni--Filmus--Lifshitz--Lindzey--Vinyals
\cite{FilmusIhringerJunta,ChiarelliHatamiSaks,Midrijanis,
FilmusJuntaThreshold,DafniFilmusLifshitzLindzeyVinyals}.  In particular, the
last authors already prove polynomial relations between degree and query
complexity on slices.  None of these inputs is claimed as new.  The
contribution here is the row-space support pooling, synchronization of
non-unique slice supports, symmetric-tree canonicalization, and the resulting
rank-sensitive communication theorems, including the bounded-alphabet product
and multislice extensions.  The remaining quantitative gap is
between the sharp $2s$ threshold with unspecified dependence, the polynomial
$144s^4$ query regime, and the sharper $s$-dependence in the $C^s$
cube-extension regime.

Concept-class evaluation is used by Feldman--Xiao to relate randomized
one-way communication to differentially private learning
\cite{FeldmanXiao}.  Mande--Sanyal--Sherif relate one-way communication and
non-adaptive decision trees for functions composed with fixed gadgets
\cite{MandeSanyalSherif}.  These works do not state a deterministic
real-rank-sensitive bounded-answer evaluation theorem.  Theorem
\ref{thm:rank-pooled-query} instead applies directly to an arbitrary Boolean
evaluation matrix and replaces the ambient coordinate cost by its real row
rank.  This targeted comparison is still not a substitute for an independent
priority search in learning theory.

The sparse-column extension uses standard constant-leverage sampling on a
transitive invariant subspace, matrix Chernoff, and the Gross--Nesme
without-replacement comparison \cite{Tropp,GrossNesme}.  Those concentration
ingredients are not new.  The claim here is the unisolvent rank-lifting step
and its combination with Theorem~\ref{thm:rank-pooled-query}, yielding one
sparse column family that works simultaneously for all eligible Boolean row
families.  The twice-degree certificate $f(f-1)$ additionally removes the
need to assume that rows are Boolean outside the sparse family; this is the
step used in Theorem~\ref{thm:quadratic-column-hls}.

Amireddy--Behera--Srinivasan--Sudan--Willumsgaard recently construct
constant-query randomized low-degree tests on balanced Boolean slices
\cite{AmireddyBeheraSrinivasanSudanWillumsgaard}.  Their goal is to
distinguish low degree from functions far from low degree using fresh random
oracle queries.  Theorems~\ref{thm:booleanity-certifying-transfer} and
\ref{thm:quadratic-column-hls} instead use one fixed minimal interpolation
set, certify exact Booleanity for every affine row simultaneously, and derive
a deterministic real-rank communication protocol.

Maximal rank and diagonal forms for the full pair-versus-$k$-set inclusion
matrix are classical \cite{GottliebIncidence,WilsonIncidence}.  Those results
imply that some square full-rank minor exists.  Theorem
\ref{thm:explicit-cross} identifies a specific base/single-swap/two-double-star
minor and proves it nonsingular by direct coefficient elimination.  The
incidence-rank fact itself is not a novelty claim; candidate new content is
the explicit sparse minor and its universal adjacent-two-value communication
application.

The entropy estimate used here is elementary.  It is compatible with the
standard discrete entropy--variance inequality obtained by adding independent
uniform noise and applying Gaussian maximal entropy
\cite{GavalakisKontoyiannis}.

Targeted searches found no prior theorem with the product hypothesis
\eqref{eq:product-family} and relative bound \eqref{eq:master-bound}, nor a
dimension-free rectangle theorem parameterized by Boolean tube radius and
integer-generator entropy.  This is not a substitute for an independent
publication-level novelty audit.

\section{Limitations and open problems}

The spectral--lattice hypotheses are sufficient, not necessary.  An
irrational low-mode space can have useful Boolean structure not captured by
integer generators.  Conversely, a rational approximation may have poor
Boolean margin or prohibitively large height.  The robust theorem isolates
the remaining inverse task: produce a projector-close rational space with
simultaneously controlled Boolean margin and arithmetic entropy, or extract a
coordinate/equality/code block directly.

The product-Johnson theorem has a full logarithmic one-way protocol because
its rows admit the exact coordinate/complement classification.  In contrast,
the general spectral--lattice theorem is presently only a one-rectangle
theorem: arbitrary recursive restrictions need not preserve its covariance,
tube-radius, or generator-entropy hypotheses.

The most valuable next directions are:
\begin{enumerate}
\item reduce the polynomial depth condition $144s^4$ toward the optimal
$2s$ junta threshold while retaining explicit dependence on $s$;
\item unequal and near-boundary product layers with sharp tube profiles;
\item random and pseudorandom thinning inside each block;
\item a relative inverse theorem producing the rational data required by
Theorem~\ref{thm:robust};
\item an external comparison with association-scheme and extremal product
theorems not indexed by the terminology used here.
\end{enumerate}

\section{Machine verification}

The accompanying reviewer repository contains a curated 18-module Lean 4.22
core, reduced from the broader historical research development.  It contains
no \texttt{sorry}, \texttt{admit}, project-defined axioms, compiled objects,
or cached dependencies.  The retained modules formalize the finite cores
needed for the headline results:
\begin{itemize}
\item two-point variance and the explicit rank-at-most-$d+1$ factorization;
\item unisolvent restriction, affine support pairing, and rank-basis support
pooling;
\item uniform-layer and product-Johnson range/kernel lemmas;
\item the short-side, low-density, and global six-coordinate sharp $2/5$
polynomials, including the unique $J(6,3)$ equality and the $k=2$ boundary;
\item the deterministic covariance-gap medium-row contradiction and the
Sylvester--Hadamard threshold.
\end{itemize}
The probabilistic sampling inputs (matrix Chernoff and VC), the projection
DPP, published junta and query theorems, and the complete protocol wrapper
remain conventional paper proofs.  The exact coverage boundary is recorded in
\nolinkurl{formal/FORMALIZATION_MANIFEST.md}.

Five assertion-based Python suites provide independent finite regression.
The current checks include 2,533,728 product-Johnson row--column pairs, 5,400
general product protocol cases, 1,000 product-multislice cases, 105,747 slice
support-intersection instances, 16,952 symmetric-tree instances, 155,155
hypergeometric parameter triples, 2,047,045 adjacent intervals, 123,750
six-coordinate certificate parameters, 9,241,290 affine evaluations, and six
Sylvester--Hadamard designs.  The sharp audit finds the unique normalized
$2/5$ equality $(n,k,|A|,a)=(6,3,3,1)$.

The repository provides one-command targets for the Lean build, all Python
checks, PDF generation, and a minimal arXiv source archive.  Historical
Fourier, affine-section, extraction, and exploratory modules are deliberately
omitted because they are not dependencies of the reviewer-facing claims.

\bibliographystyle{alpha}
\bibliography{references}

\end{document}
